% ============================================================
%  SAMURAI — Математика, Физика и Дифференциальные системы
%  Главный файл документа (собирает все главы)
%  Компиляция: pdflatex main.tex  (два прохода для оглавления)
% ============================================================
\documentclass[12pt, a4paper, openany]{book}

% ── Кодировка и язык ─────────────────────────────────────────────
\usepackage[utf8]{inputenc}
\usepackage[T2A]{fontenc}
\usepackage[russian, english]{babel}

% ── Геометрия страницы ───────────────────────────────────────────
\usepackage[
  top=25mm, bottom=25mm,
  left=30mm, right=20mm
]{geometry}

% ── Математика ───────────────────────────────────────────────────
\usepackage{amsmath, amssymb, amsthm}
\usepackage{mathtools}
\usepackage{bm}           % жирные математические символы

% ── Графика и цвет ───────────────────────────────────────────────
\usepackage{xcolor}
\usepackage{graphicx}
\usepackage{tikz}
\usetikzlibrary{arrows.meta, positioning, shapes.geometric, calc}

% ── Код ──────────────────────────────────────────────────────────
\usepackage{listings}
\usepackage{courier}

% Цветовая схема для кода
\definecolor{codebg}{RGB}{245,247,250}
\definecolor{codeframe}{RGB}{180,190,210}
\definecolor{keyword}{RGB}{0,100,200}
\definecolor{comment}{RGB}{100,130,100}
\definecolor{string}{RGB}{163,21,21}
\definecolor{number}{RGB}{9,134,88}

% Стиль Python
\lstdefinestyle{pythonstyle}{
  language=Python,
  backgroundcolor=\color{codebg},
  frame=single,
  rulecolor=\color{codeframe},
  framesep=4pt,
  basicstyle=\footnotesize\ttfamily,
  keywordstyle=\color{keyword}\bfseries,
  commentstyle=\color{comment}\itshape,
  stringstyle=\color{string},
  numberstyle=\scriptsize\ttfamily\color{gray},
  numbers=left,
  numbersep=10pt,
  xleftmargin=2.5em,
  stepnumber=1,
  showstringspaces=false,
  breaklines=true,
  breakatwhitespace=false,
  tabsize=4,
  captionpos=b,
  morekeywords={True, False, None, self},
}

% Стиль C++
\lstdefinestyle{cppstyle}{
  language=C++,
  backgroundcolor=\color{codebg},
  frame=single,
  rulecolor=\color{codeframe},
  framesep=4pt,
  basicstyle=\footnotesize\ttfamily,
  keywordstyle=\color{keyword}\bfseries,
  commentstyle=\color{comment}\itshape,
  stringstyle=\color{string},
  numberstyle=\scriptsize\ttfamily\color{gray},
  numbers=left,
  numbersep=10pt,
  xleftmargin=2.5em,
  stepnumber=1,
  showstringspaces=false,
  breaklines=true,
  tabsize=4,
  captionpos=b,
}

\lstset{style=pythonstyle}

% ── Таблицы и плавающие объекты ──────────────────────────────────
\usepackage{booktabs}
\usepackage{array}
\usepackage{longtable}
\usepackage{float}

% ── Гиперссылки ──────────────────────────────────────────────────
\usepackage[
  colorlinks=true,
  linkcolor=blue!70!black,
  citecolor=green!60!black,
  urlcolor=blue!80!black,
  bookmarks=true,
  bookmarksnumbered=true
]{hyperref}

% ── Колонтитулы ──────────────────────────────────────────────────
\usepackage{fancyhdr}
\pagestyle{fancy}
\fancyhf{}
\fancyhead[LE,RO]{\small\thepage}
\fancyhead[LO]{\small\nouppercase{\rightmark}}
\fancyhead[RE]{\small\nouppercase{\leftmark}}
\renewcommand{\headrulewidth}{0.4pt}

% ── Параграфы ────────────────────────────────────────────────────
\usepackage{parskip}
\setlength{\parindent}{0pt}

% ── Окружения ────────────────────────────────────────────────────
\usepackage{tcolorbox}
\tcbuselibrary{skins, breakable}

\newtcolorbox{formulabox}[1][]{
  enhanced,
  breakable,
  colback=blue!3!white,
  colframe=blue!50!black,
  fonttitle=\bfseries,
  title=Формула,
  #1
}

\newtcolorbox{notebox}[1][]{
  enhanced,
  colback=orange!8!white,
  colframe=orange!70!black,
  fonttitle=\bfseries,
  title=Примечание,
  #1
}

\newtcolorbox{filebox}[1][]{
  enhanced,
  colback=gray!6!white,
  colframe=gray!50!black,
  fonttitle=\bfseries\ttfamily,
  #1
}

% ── Нумерация уравнений по главам ────────────────────────────────
\numberwithin{equation}{chapter}

% ── Макросы ──────────────────────────────────────────────────────
\newcommand{\filepath}[1]{\texttt{\small\color{gray!80!black}#1}}
\newcommand{\vect}[1]{\bm{#1}}
\newcommand{\mat}[1]{\mathbf{#1}}
\newcommand{\T}{^{\mathsf{T}}}

% ============================================================
\begin{document}

% ── Титульная страница ───────────────────────────────────────────
\begin{titlepage}
\centering
\vspace*{2cm}
{\Huge\bfseries SAMURAI Robot}\\[0.5cm]
{\Large\bfseries Математика, Физика и Дифференциальные системы}\\[1cm]
\rule{\textwidth}{1pt}\\[0.5cm]
{\large Технический справочник по математическому ядру проекта}\\[0.5cm]
\rule{\textwidth}{1pt}\\[2cm]

\begin{tabular}{ll}
\textbf{Фильтры состояния:} & Extended Kalman Filter (2D velocity, IMU yaw)\\
\textbf{Оценка позиции:} & Акселерометр + одометрия, слияние через EKF\\
\textbf{Управление:} & Pure Pursuit, пропорциональный регулятор\\
\textbf{Картография:} & SLAM на log-odds, трассировка лучей\\
\textbf{Симулятор:} & Кинематика, пересечение луч--сфера, перспектива\\
\textbf{Алгоритмы:} & A*, Bresenham, проекции, кватернионы\\
\end{tabular}

\vfill
{\large\today}
\end{titlepage}

% ── Оглавление ───────────────────────────────────────────────────
\tableofcontents
\newpage

% ── Введение ─────────────────────────────────────────────────────
\chapter*{Введение}
\addcontentsline{toc}{chapter}{Введение}

Данный документ является техническим справочником по всем математическим,
физическим и дифференциальным системам, реализованным в проекте \textbf{SAMURAI} ---
автономного наземного робота на гусеничном шасси.

Документ построен \textbf{модульно}: каждая глава соответствует отдельному
программному модулю и содержит:
\begin{itemize}
  \item Ссылку на исходный файл с указанием строк кода
  \item Математическое описание алгоритма
  \item Точный фрагмент кода с пояснениями
  \item Физический смысл каждого уравнения
\end{itemize}

\textbf{Структура проекта (математические компоненты):}
\begin{center}
\begin{tabular}{lll}
\toprule
\textbf{Файл} & \textbf{Язык} & \textbf{Содержание}\\
\midrule
\texttt{pi\_nodes/filters/velocity\_ekf.py} & Python & 2D Velocity EKF\\
\texttt{pi\_nodes/filters/ekf\_imu.py} & Python & IMU yaw 1D Kalman\\
\texttt{pi\_nodes/filters/accel\_position.py} & Python & Слияние датчиков\\
\texttt{pi\_nodes/nodes/path\_recorder\_node.py} & Python & Pure Pursuit\\
\texttt{pi\_nodes/nodes/slam\_map\_node.py} & Python & SLAM log-odds\\
\texttt{pi\_nodes/nodes/imu\_node.py} & Python & IMU калибровка + EMA\\
\texttt{pi\_nodes/nodes/fallback\_nav\_node.py} & Python & Реактивная навигация\\
\texttt{compute\_node/simulator.py} & Python & Физика + A* + Bresenham\\
\texttt{compute\_node/dashboard\_node.py} & Python & Кватернионы\\
\texttt{compute\_node/geofence\_map\_publisher.py} & Python & Point-in-Polygon\\
\texttt{compute\_node/depth\_to\_scan\_node.py} & Python & УЗ $\to$ LaserScan\\
\texttt{ros\_ws/.../follow\_me\_node.py} & Python & Pinhole-камера + P-рег.\\
\texttt{ros\_ws/.../patrol\_node.py} & Python & Навигация + кватернионы\\
\texttt{ros\_ws/.../motor\_node.cpp} & C++ & Одометрия + ШИМ-микшер\\
\texttt{ros\_ws/.../imu\_node.cpp} & C++ & I2C IMU (оптимизация)\\
\bottomrule
\end{tabular}
\end{center}

% ── Главы ────────────────────────────────────────────────────────
% ============================================================
%  Глава 0 — Теоретические основы теории управления
%  Справочная глава: модели, линеаризация, устойчивость
% ============================================================
\chapter{Теоретические основы систем управления}
\label{ch:theory}

Данная глава содержит математический аппарат теории автоматического
управления (ТАУ), на котором базируются все алгоритмы проекта SAMURAI.

% ═══════════════════════════════════════════════════════════════
\section{Модель в пространстве состояний}
\label{sec:state_space}

Любой объект управления может быть описан системой дифференциальных
уравнений первого порядка --- \textbf{уравнениями состояния}:

\begin{equation}
  \boxed{
  \dot{\vect{x}}(t) = \mat{A}\,\vect{x}(t) + \mat{B}\,\vect{u}(t), \qquad
  \vect{y}(t) = \mat{C}\,\vect{x}(t) + \mat{D}\,\vect{u}(t)
  }
  \label{eq:state_space}
\end{equation}

где:
\begin{itemize}
  \item $\vect{x}(t) = (x_1, x_2, \ldots, x_n)\T \in \mathbb{R}^n$ --- \textbf{вектор состояния}
  \item $\vect{u}(t) = (u_1, \ldots, u_r)\T \in \mathbb{R}^r$ --- вектор управляющих воздействий
  \item $\vect{y}(t) = (y_1, \ldots, y_m)\T \in \mathbb{R}^m$ --- вектор выходных переменных
  \item $\mat{A} \in \mathbb{R}^{n \times n}$ --- матрица системы (динамики)
  \item $\mat{B} \in \mathbb{R}^{n \times r}$ --- матрица управления
  \item $\mat{C} \in \mathbb{R}^{m \times n}$ --- матрица наблюдения
  \item $\mat{D} \in \mathbb{R}^{m \times r}$ --- матрица прямого прохода
\end{itemize}

\begin{notebox}
Структурная схема модели в пространстве состояний содержит
интегратор ($\int$) для перехода от $\dot{\vect{x}}$ к $\vect{x}$,
блоки умножения на матрицы $\mat{A}$, $\mat{B}$, $\mat{C}$, $\mat{D}$
и сумматоры.
\end{notebox}

\subsection{Применение в проекте SAMURAI}

Каждый фильтр Калмана в проекте является реализацией модели
в пространстве состояний:

\begin{center}
\begin{tabular}{lll}
\toprule
\textbf{Модуль} & \textbf{Вектор состояния $\vect{x}$} & \textbf{Размерность $n$}\\
\midrule
Velocity EKF (гл.~\ref{ch:velocity_ekf}) & $(v_x, v_y)\T$ & 2\\
IMU EKF (гл.~\ref{ch:imu_ekf}) & $(\psi, b_z)\T$ & 2\\
AccelPosition (гл.~\ref{ch:accel_position}) & $(x, y, v_x, v_y)\T$ & 4\\
Одометрия (гл.~\ref{ch:odometry}) & $(x, y, \theta)\T$ & 3\\
\bottomrule
\end{tabular}
\end{center}

% ═══════════════════════════════════════════════════════════════
\section{Линеаризация нелинейных систем}
\label{sec:linearization}

Реальные объекты, включая наземного робота, описываются
\textbf{нелинейными} уравнениями:

\begin{equation}
  \dot{\vect{x}}(t) = \vect{f}(\vect{x}, \vect{u}), \qquad
  \vect{y}(t) = \vect{g}(\vect{x}, \vect{u})
  \label{eq:nonlinear_system}
\end{equation}

\subsection{Пример: транспортная тележка}

Модель движения мобильного робота (транспортной тележки) описывается
системой шести дифференциальных уравнений:

\begin{align}
  \dot{\theta} &= -\frac{1}{T_1}\theta + \frac{1}{T_1}u_1 \label{eq:cart_theta}\\
  \dot{f} &= -\frac{1}{T_2}f + \frac{1}{T_2}u_2 \label{eq:cart_force}\\
  \dot{v} &= -\frac{\alpha}{m}v + \frac{1}{m}f \label{eq:cart_vel}\\
  \dot{x} &= v\cos(\psi + \theta) \label{eq:cart_x}\\
  \dot{\psi} &= \frac{1}{L}v\sin(\theta) \label{eq:cart_psi}\\
  \dot{z} &= -v\sin(\psi + \theta) \label{eq:cart_z}
\end{align}

где $\theta$ --- угол поворота колёс, $\psi$ --- угол рыскания,
$f$ --- движущая сила, $v$ --- скорость, $x,z$ --- координаты;
$T_1, T_2, \alpha, m, L$ --- параметры объекта.

Уравнения~\eqref{eq:cart_x}--\eqref{eq:cart_z} \textbf{нелинейны}
(содержат $\sin$ и $\cos$), что делает систему нелинейной.

\begin{notebox}
Модель робота SAMURAI использует ту же структуру:
уравнения одометрии (гл.~\ref{ch:odometry}) $\dot{x} = v\cos\theta$,
$\dot{y} = v\sin\theta$ нелинейны из-за тригонометрических функций
от угла курса $\theta$.
\end{notebox}

\subsection{Метод линеаризации по Тейлору (матрица Якоби)}
\label{sec:jacobian}

Линеаризация проводится в окрестности \textbf{опорного режима}
$(\bar{\vect{x}}, \bar{\vect{u}})$, в котором $\vect{f}(\bar{\vect{x}}, \bar{\vect{u}}) = \vect{0}$.

Разложение в ряд Тейлора с отбрасыванием членов высших порядков:

\begin{equation}
  \boxed{
  \Delta\dot{\vect{x}} = \mat{A}\,\Delta\vect{x} + \mat{B}\,\Delta\vect{u}, \qquad
  \Delta\vect{y} = \mat{C}\,\Delta\vect{x} + \mat{D}\,\Delta\vect{u}
  }
  \label{eq:linearized}
\end{equation}

где матрицы --- \textbf{якобианы}, вычисленные в опорной точке:

\begin{equation}
  \mat{A} = \left.\frac{\partial \vect{f}}{\partial \vect{x}}\right|_{\bar{\vect{x}},\bar{\vect{u}}}, \quad
  \mat{B} = \left.\frac{\partial \vect{f}}{\partial \vect{u}}\right|_{\bar{\vect{x}},\bar{\vect{u}}}, \quad
  \mat{C} = \left.\frac{\partial \vect{g}}{\partial \vect{x}}\right|_{\bar{\vect{x}},\bar{\vect{u}}}, \quad
  \mat{D} = \left.\frac{\partial \vect{g}}{\partial \vect{u}}\right|_{\bar{\vect{x}},\bar{\vect{u}}}
  \label{eq:jacobians}
\end{equation}

\subsection{Пример: якобиан для транспортной тележки}

Для опорного режима: $v(t) = v^0$, $\theta(t) = 0$, $\psi(t) = 0$:

\begin{equation}
  \mat{A} = \begin{pmatrix}
    -1/T_1 & 0 & 0 & 0 & 0 & 0\\
    0 & -1/T_2 & 0 & 0 & 0 & 0\\
    0 & 1/m & -\alpha/m & 0 & 0 & 0\\
    0 & 0 & 1 & 0 & 0 & 0\\
    v^0/L & 0 & 0 & 0 & 0 & 0\\
    -v^0 & 0 & 0 & 0 & -v^0 & 0
  \end{pmatrix}
  \label{eq:cart_A}
\end{equation}

\subsection{Связь с EKF в проекте SAMURAI}

\textbf{Расширенный фильтр Калмана (EKF)} --- это именно
\textbf{линеаризация на каждом шаге}. Матрица перехода $\mat{F}$
в фильтре Калмана (гл.~\ref{ch:imu_ekf}, формула~\eqref{eq:F_matrix})
--- это якобиан нелинейной модели:

\begin{equation}
  \mat{F}_k = \left.\frac{\partial \vect{f}}{\partial \vect{x}}\right|_{\vect{x}_{k-1}}
  = \begin{pmatrix} 1 & -\Delta t \\ 0 & 1 \end{pmatrix}
  \label{eq:ekf_jacobian}
\end{equation}

Здесь нелинейная модель $\psi_k = \psi_{k-1} + (g_z - b_z)\Delta t$
линеаризуется относительно текущего состояния $(\psi_{k-1}, b_{z,k-1})$,
а якобиан $\mat{F}$ содержит производные $\partial\dot{\psi}/\partial\psi = 1$,
$\partial\dot{\psi}/\partial b_z = -\Delta t$.

% ═══════════════════════════════════════════════════════════════
\section{Передаточная функция}
\label{sec:transfer_func}

Для линейной стационарной системы~\eqref{eq:state_space}
\textbf{передаточная функция} определяется через преобразование Лапласа:

\begin{equation}
  \boxed{W(p) = \mat{C}(p\mat{E}_n - \mat{A})^{-1}\mat{B} + \mat{D}}
  \label{eq:transfer_func}
\end{equation}

где $p = \frac{d}{dt}$ --- оператор дифференцирования, $\mat{E}_n$ --- единичная матрица.

\subsection{Характеристический полином и резольвента}

Резольвента матрицы $\mat{A}$:

\begin{equation}
  (p\mat{E}_n - \mat{A})^{-1} = \frac{1}{\chi_n(p)}\widetilde{\mat{B}}(p)
  \label{eq:resolvent}
\end{equation}

где $\chi_n(p) = \det(p\mat{E}_n - \mat{A}) = p^n + a_1 p^{n-1} + \ldots + a_n$
--- характеристический полином, а $\widetilde{\mat{B}}(p)$ --- присоединённая матрица.

\subsection{Алгоритм Фаддеева--Леверье}
\label{sec:faddeev}

Коэффициенты характеристического полинома и матрицы
$\widetilde{\mat{B}}_k$ вычисляются рекуррентно:

\begin{equation}
  \boxed{
  \begin{aligned}
    \widetilde{\mat{B}}_1 &= \mat{E}_n, \quad & a_1 &= -\mathrm{tr}(\mat{A}\widetilde{\mat{B}}_1)\\
    \widetilde{\mat{B}}_k &= \mat{A}\widetilde{\mat{B}}_{k-1} + a_{k-1}\mat{E}_n, \quad
    & a_k &= -\frac{1}{k}\mathrm{tr}(\mat{A}\widetilde{\mat{B}}_k)\\
    & & \multicolumn{2}{l}{k = 2, \ldots, n}
  \end{aligned}
  }
  \label{eq:faddeev}
\end{equation}

Проверка: $\mat{A}\widetilde{\mat{B}}_n + a_n\mat{E}_n = \mat{0}$
(теорема Гамильтона--Кели).

\subsection{Уравнение <<вход--выход>>}

Передаточная функция связывает выход $y$ с входом $u$ через
дифференциальное уравнение:

\begin{equation}
  p^n y + a_1 p^{n-1} y + \ldots + a_n y = b_0 p^n u + b_1 p^{n-1} u + \ldots + b_n u
  \label{eq:io_equation}
\end{equation}

\begin{notebox}
В проекте SAMURAI фильтры Калмана работают в \textbf{дискретном}
пространстве состояний:
$\vect{x}_{k+1} = \mat{F}\,\vect{x}_k + \mat{B}\,\vect{u}_k$,
что является дискретным аналогом непрерывной модели~\eqref{eq:state_space}.
Переход к дискретной модели:
$\mat{A}_d = e^{\mat{A}T}$, $\mat{B}_d = \int_0^T e^{\mat{A}\tau}d\tau\,\mat{B}$.
\end{notebox}

% ═══════════════════════════════════════════════════════════════
\section{Устойчивость систем управления}
\label{sec:stability}

\subsection{Определение устойчивости}

Система~\eqref{eq:state_space} \textbf{асимптотически устойчива},
если решение $\vect{x}(t) \to \vect{0}$ при $t \to \infty$
для любых начальных условий $\vect{x}(0)$.

\subsection{Критерий собственных значений}

Система устойчива тогда и только тогда, когда все собственные числа
матрицы $\mat{A}$ имеют отрицательные вещественные части:

\begin{equation}
  \boxed{\mathrm{Re}(\lambda_i) < 0, \quad \forall\, i = 1, \ldots, n}
  \label{eq:stability_eigenvalues}
\end{equation}

где $\lambda_i$ --- корни характеристического уравнения $\det(\lambda\mat{E} - \mat{A}) = 0$.

\subsection{Критерий Гурвица}

Для характеристического полинома $\chi(p) = p^n + a_1 p^{n-1} + \ldots + a_n$
формируется матрица Гурвица:

\begin{equation}
  \mat{H} = \begin{pmatrix}
    a_1 & a_3 & a_5 & \cdots\\
    1   & a_2 & a_4 & \cdots\\
    0   & a_1 & a_3 & \cdots\\
    0   & 1   & a_2 & \cdots\\
    \vdots & & & \ddots
  \end{pmatrix}
  \label{eq:hurwitz}
\end{equation}

\textbf{Условие устойчивости}: все угловые миноры матрицы $\mat{H}$
положительны:

\begin{equation}
  \Delta_1 = a_1 > 0, \quad
  \Delta_2 = a_1 a_2 - a_3 > 0, \quad
  \ldots, \quad \Delta_n > 0
  \label{eq:hurwitz_cond}
\end{equation}

\subsection{Решение уравнения состояния}

Аналитическое решение линейной системы (формула Коши):

\begin{equation}
  \vect{x}(t) = e^{\mat{A}(t - t_0)}\vect{x}(t_0) + \int_{t_0}^{t} e^{\mat{A}(t - \tau)}\mat{B}\,\vect{u}(\tau)\,d\tau
  \label{eq:cauchy}
\end{equation}

где $e^{\mat{A}t}$ --- матричная экспонента. Первое слагаемое ---
свободное движение (от начальных условий), второе ---
вынужденное движение (от входного воздействия).

\subsection{Жорданова форма и матричная экспонента}

Через замену $\vect{x} = \mat{S}\vect{z}$ система приводится к жордановой форме:

\begin{equation}
  \dot{\vect{z}} = \mat{J}\vect{z} + \mat{S}^{-1}\mat{B}\vect{u}, \quad
  \mat{J} = \mat{S}^{-1}\mat{A}\mat{S}
  \label{eq:jordan}
\end{equation}

Для диагональной жордановой формы экспонента:

\begin{equation}
  e^{\mat{J}t} = \mathrm{diag}\!\left(e^{\lambda_1 t}, e^{\lambda_2 t}, \ldots, e^{\lambda_n t}\right)
  \label{eq:matrix_exp}
\end{equation}

Отсюда видно: при $\mathrm{Re}(\lambda_i) < 0$ каждая мода затухает
экспоненциально, что обеспечивает устойчивость.

% ═══════════════════════════════════════════════════════════════
\section{Дискретные системы и фильтр Калмана}
\label{sec:discrete_kalman}

\subsection{Дискретная модель в пространстве состояний}

\begin{equation}
  \vect{x}_{k+1} = \mat{A}_d\,\vect{x}_k + \mat{B}_d\,\vect{u}_k, \qquad
  \vect{y}_k = \mat{C}_d\,\vect{x}_k + \mat{D}_d\,\vect{u}_k
  \label{eq:discrete_ss}
\end{equation}

где $\mat{A}_d = e^{\mat{A}T}$, $T$ --- период дискретизации.

\subsection{Общие уравнения фильтра Калмана}

Фильтр Калмана решает задачу оптимальной оценки состояния
стохастической системы:

\begin{equation}
  \vect{x}_{k} = \mat{F}\,\vect{x}_{k-1} + \mat{B}\,\vect{u}_{k-1} + \vect{w}_{k-1}, \qquad
  \vect{z}_k = \mat{H}\,\vect{x}_k + \vect{v}_k
  \label{eq:kalman_model}
\end{equation}

где $\vect{w} \sim \mathcal{N}(\vect{0}, \mat{Q})$ --- процессный шум,
$\vect{v} \sim \mathcal{N}(\vect{0}, \mat{R})$ --- шум измерений.

\textbf{Шаг предсказания (Predict):}

\begin{align}
  \hat{\vect{x}}_k^- &= \mat{F}\,\hat{\vect{x}}_{k-1} + \mat{B}\,\vect{u}_{k-1}
  \label{eq:kf_predict_x}\\
  \mat{P}_k^- &= \mat{F}\,\mat{P}_{k-1}\,\mat{F}\T + \mat{Q}
  \label{eq:kf_predict_P}
\end{align}

\textbf{Шаг коррекции (Update):}

\begin{align}
  \vect{y}_k &= \vect{z}_k - \mat{H}\,\hat{\vect{x}}_k^- && \text{(инновация)}
  \label{eq:kf_innovation}\\
  \mat{S}_k &= \mat{H}\,\mat{P}_k^-\,\mat{H}\T + \mat{R} && \text{(ковариация инновации)}
  \label{eq:kf_S}\\
  \mat{K}_k &= \mat{P}_k^-\,\mat{H}\T\,\mat{S}_k^{-1} && \text{(усиление Калмана)}
  \label{eq:kf_gain}\\
  \hat{\vect{x}}_k &= \hat{\vect{x}}_k^- + \mat{K}_k\,\vect{y}_k && \text{(обновление состояния)}
  \label{eq:kf_update_x}\\
  \mat{P}_k &= (\mat{I} - \mat{K}_k\mat{H})\,\mat{P}_k^- && \text{(обновление ковариации)}
  \label{eq:kf_update_P}
\end{align}

\begin{formulabox}
Фильтр Калмана \textbf{минимизирует} дисперсию ошибки оценки
$E[\|\vect{x} - \hat{\vect{x}}\|^2]$. Усиление $\mat{K}$ автоматически
взвешивает доверие между предсказанием (модель) и измерением (датчик):
\begin{itemize}
  \item $\mat{R} \ll \mat{P}^-$: $\mat{K} \to \mat{H}^{-1}$ --- доверяем датчику
  \item $\mat{R} \gg \mat{P}^-$: $\mat{K} \to \mat{0}$ --- доверяем модели
\end{itemize}
\end{formulabox}

\subsection{Расширенный фильтр Калмана (EKF)}

Для нелинейной системы $\vect{x}_k = \vect{f}(\vect{x}_{k-1}, \vect{u}_{k-1})$,
$\vect{z}_k = \vect{h}(\vect{x}_k)$ используется \textbf{линеаризация
на каждом шаге}:

\begin{equation}
  \mat{F}_k = \left.\frac{\partial \vect{f}}{\partial \vect{x}}\right|_{\hat{\vect{x}}_{k-1}}, \qquad
  \mat{H}_k = \left.\frac{\partial \vect{h}}{\partial \vect{x}}\right|_{\hat{\vect{x}}_k^-}
  \label{eq:ekf_linearize}
\end{equation}

Затем якобианы $\mat{F}_k$, $\mat{H}_k$ подставляются в стандартные
уравнения фильтра Калмана~\eqref{eq:kf_predict_P}--\eqref{eq:kf_update_P}.

% ═══════════════════════════════════════════════════════════════
\section{Модель электродвигателя постоянного тока}
\label{sec:dc_motor}

Двигатели робота SAMURAI управляются ШИМ-сигналами через драйвер TB6612.
Электромеханические процессы в двигателе постоянного тока описываются
дифференциальным уравнением второго порядка:

\begin{equation}
  \boxed{T_a^2\ddot{\omega} + T_y\dot{\omega} + \omega = k\,u}
  \label{eq:dc_motor}
\end{equation}

где:
\begin{itemize}
  \item $\omega$ --- угловая скорость вала двигателя
  \item $u$ --- управляющее воздействие (напряжение / ШИМ)
  \item $T_a$ --- электромеханическая постоянная времени
  \item $T_y$ --- электрическая постоянная времени
  \item $k$ --- коэффициент передачи
\end{itemize}

\subsection{Пространство состояний для двигателя}

Вводя $x_1 = \omega$, $x_2 = \dot{\omega}$:

\begin{equation}
  \begin{pmatrix} \dot{x}_1 \\ \dot{x}_2 \end{pmatrix}
  = \begin{pmatrix} 0 & 1 \\ -1/T_a^2 & -T_y/T_a^2 \end{pmatrix}
  \begin{pmatrix} x_1 \\ x_2 \end{pmatrix}
  + \begin{pmatrix} 0 \\ k/T_a^2 \end{pmatrix} u
  \label{eq:motor_ss}
\end{equation}

Характеристический полином:

\begin{equation}
  \chi(\lambda) = \lambda^2 + \frac{T_y}{T_a^2}\lambda + \frac{1}{T_a^2}
  \label{eq:motor_char}
\end{equation}

\textbf{Условие устойчивости} (по Гурвицу): $T_y > 0$, $T_a > 0$ ---
выполняется для всех реальных двигателей.

\subsection{Передаточная функция двигателя}

\begin{equation}
  W(p) = \frac{k}{T_a^2 p^2 + T_y p + 1}
  \label{eq:motor_tf}
\end{equation}

При $T_a \ll T_y$ (типично для малых моторов) двигатель приближённо
описывается звеном первого порядка:

\begin{equation}
  W(p) \approx \frac{k}{T_y p + 1}
  \label{eq:motor_tf_approx}
\end{equation}

\begin{notebox}
В проекте SAMURAI ШИМ-управление (гл.~\ref{ch:pwm_motor})
работает с нормированным throttle $t \in [-1, 1]$, что соответствует
$u/u_{\max}$. Драйвер TB6612 в режиме Slow Decay обеспечивает
линейную зависимость $\omega \propto t$ в установившемся режиме.
\end{notebox}

% ═══════════════════════════════════════════════════════════════
\section{Пропорциональный регулятор (P-регулятор)}
\label{sec:p_controller}

\subsection{Закон управления}

\begin{equation}
  u(t) = K_P \cdot e(t), \qquad e(t) = r(t) - y(t)
  \label{eq:p_controller}
\end{equation}

где $r(t)$ --- задание, $y(t)$ --- выход, $K_P$ --- коэффициент усиления.

\subsection{Передаточная функция замкнутой системы}

Для объекта с передаточной функцией $W_{\text{obj}}(p)$
и P-регулятора с коэффициентом $K_P$:

\begin{equation}
  W_{\text{зс}}(p) = \frac{K_P \cdot W_{\text{obj}}(p)}{1 + K_P \cdot W_{\text{obj}}(p)}
  \label{eq:closed_loop}
\end{equation}

\subsection{Статическая ошибка}

P-регулятор не может полностью устранить установившуюся ошибку
при ступенчатом воздействии:

\begin{equation}
  e_{\text{уст}} = \frac{1}{1 + K_P \cdot k_{\text{obj}}}
  \label{eq:static_error}
\end{equation}

Увеличение $K_P$ снижает ошибку, но может привести к
потере устойчивости.

\subsection{Применение P-регуляторов в проекте}

\begin{center}
\begin{tabular}{llll}
\toprule
\textbf{Модуль} & \textbf{Управляемая переменная} & $K_P$ & \textbf{Ограничение}\\
\midrule
Pure Pursuit (гл.~\ref{ch:pure_pursuit}) & Угловая скорость $\omega$ & $2.5$ & $|\omega| \leq 0.8$\\
Follow-Me (гл.~\ref{ch:follow_me}) & Угловая скорость $\omega$ & $0.8$ & ---\\
Follow-Me (гл.~\ref{ch:follow_me}) & Линейная скорость $v$ & $0.3$ & $v \leq 0.20$\\
Patrol (гл.~\ref{ch:patrol}) & Целевая позиция & --- & $d < 0.20\,\text{м}$\\
\bottomrule
\end{tabular}
\end{center}

% ═══════════════════════════════════════════════════════════════
\section{Переходная характеристика}
\label{sec:transient}

Переходная характеристика $h(t)$ --- реакция системы на единичное
ступенчатое воздействие $u(t) = 1(t)$ при нулевых начальных условиях:

\begin{equation}
  h(t) = \mat{C}\int_0^t e^{\mat{A}\tau}\,d\tau\,\mat{B} + \mat{D}
  \label{eq:step_response}
\end{equation}

При отсутствии нулевых собственных значений у $\mat{A}$:

\begin{equation}
  h(t) = \mat{C}\left(e^{\mat{A}t} - \mat{E}\right)\mat{A}^{-1}\mat{B} + \mat{D}
  \label{eq:step_response2}
\end{equation}

Для системы второго порядка с $\omega_n$ (собственная частота)
и $\zeta$ (коэффициент демпфирования):

\begin{equation}
  h(t) = 1 - \frac{e^{-\zeta\omega_n t}}{\sqrt{1-\zeta^2}}
  \sin\!\left(\omega_n\sqrt{1-\zeta^2}\,t + \arccos\zeta\right)
  \label{eq:2nd_order_response}
\end{equation}

\begin{notebox}
Косинусное масштабирование скорости в Pure Pursuit
(гл.~\ref{ch:pure_pursuit}, формула~\eqref{eq:cos_factor})
обеспечивает плавную переходную характеристику без перерегулирования
при развороте робота.
\end{notebox}

% ═══════════════════════════════════════════════════════════════
\section{Управляемость и наблюдаемость}
\label{sec:controllability}

\subsection{Матрица управляемости}

\begin{equation}
  \mat{S}_v = \bigl[\mat{B},\; \mat{A}\mat{B},\; \mat{A}^2\mat{B},\; \ldots,\; \mat{A}^{n-1}\mat{B}\bigr]
  \label{eq:controllability}
\end{equation}

Система полностью управляема $\Leftrightarrow$ $\mathrm{rank}(\mat{S}_v) = n$.

\subsection{Матрица наблюдаемости}

\begin{equation}
  \mat{S}_o = \begin{pmatrix} \mat{C} \\ \mat{C}\mat{A} \\ \vdots \\ \mat{C}\mat{A}^{n-1} \end{pmatrix}
  \label{eq:observability}
\end{equation}

Система полностью наблюдаема $\Leftrightarrow$ $\mathrm{rank}(\mat{S}_o) = n$.

\begin{notebox}
Фильтр Калмана может оценивать состояние только
\textbf{наблюдаемой} системы. EKF для угла курса
(гл.~\ref{ch:imu_ekf}) наблюдаем, так как гироскоп
измеряет $\omega_z = \dot{\psi} - b_z$, что позволяет
оценить и $\psi$, и $b_z$.
\end{notebox}

% ═══════════════════════════════════════════════════════════════
\section{Сводная таблица: теория $\to$ реализация}

\begin{center}
\begin{tabular}{lll}
\toprule
\textbf{Теоретический концепт} & \textbf{Глава} & \textbf{Реализация}\\
\midrule
$\dot{\vect{x}} = \mat{A}\vect{x} + \mat{B}\vect{u}$ & \ref{ch:velocity_ekf}, \ref{ch:imu_ekf}
  & Predict-шаг EKF\\
Якобиан $\mat{F} = \partial\vect{f}/\partial\vect{x}$ & \ref{ch:imu_ekf}
  & Матрица перехода $\mat{F}$\\
$\mat{K} = \mat{P}\mat{H}\T\mat{S}^{-1}$ & \ref{ch:velocity_ekf}
  & Усиление Калмана\\
$u = K_P \cdot e$ & \ref{ch:pure_pursuit}, \ref{ch:follow_me}
  & Угловой P-регулятор\\
$T_a^2\ddot{\omega} + T_y\dot{\omega} = ku$ & \ref{ch:pwm_motor}
  & ШИМ-управление мотором\\
$\mathrm{Re}(\lambda_i) < 0$ & \ref{ch:imu_ekf}
  & Устойчивость EKF\\
$e^{\mat{A}t}$ & \ref{ch:accel_position}
  & Интегрирование позиции\\
$\vect{x}_{k+1} = \mat{A}_d\vect{x}_k$ & \ref{ch:odometry}
  & Дискретная одометрия\\
\bottomrule
\end{tabular}
\end{center}

% ============================================================
%  Глава 1 — 2D Velocity Extended Kalman Filter
%  Файл: pi_nodes/filters/velocity_ekf.py
% ============================================================
\chapter{2D Velocity Extended Kalman Filter}
\label{ch:velocity_ekf}

\begin{filebox}[title={Файл исходного кода}]
\filepath{pi\_nodes/filters/velocity\_ekf.py} \quad (251 строка, Python)
\end{filebox}

\section{Постановка задачи}

Для точного позиционирования гусеничного робота необходимо объединить два
источника скорости:

\begin{itemize}
  \item \textbf{Колёсная одометрия} --- стабильна, не имеет дрейфа, но не учитывает
    проскальзывание гусениц на неровном покрытии.
  \item \textbf{Акселерометр (однократное интегрирование)} --- улавливает реальное
    скольжение, но накапливает ошибку интегрирования (дрейф).
\end{itemize}

Оба источника объединяются \textbf{фильтром Калмана}, который автоматически
взвешивает вклад каждого по его статистическим шумовым характеристикам.

\section{Вектор состояния и модель}

\subsection{Вектор состояния}

Состояние фильтра --- скорость робота в \emph{мировой} системе координат:

\begin{equation}
  \vect{x} = \begin{pmatrix} v_x \\ v_y \end{pmatrix} \in \mathbb{R}^2
  \label{eq:ekf_state}
\end{equation}

Позиция получается последовательным интегрированием фильтрованной скорости.

\subsection{Матрица ковариации}

Поскольку numpy не используется, матрица $2\times 2$ хранится в четырёх скалярах:

\begin{equation}
  \mat{P} = \begin{pmatrix} p_{00} & p_{01} \\ p_{10} & p_{11} \end{pmatrix}
  \label{eq:cov_matrix}
\end{equation}

Начальное значение: $\mat{P} = 0.01 \cdot \mat{I}$.

\section{Шаг предсказания (Prediction)}
\label{sec:ekf_predict}

\begin{filebox}[title={\filepath{velocity\_ekf.py}, строки 90--120 --- predict\_wheel()}]
\begin{lstlisting}[style=pythonstyle, numbers=left, firstnumber=104]
# Wheel-predicted velocity in world frame
vx_wheel = vx_cmd * cos(theta)
vy_wheel = vx_cmd * sin(theta)

# Prediction: state transitions toward wheel command
alpha = 0.5
self.vx = (1 - alpha) * self.vx + alpha * vx_wheel
self.vy = (1 - alpha) * self.vy + alpha * vy_wheel

# Process noise: covariance grows
self._p00 += self.q_wheel * dt
self._p11 += self.q_wheel * dt

# Integrate position using current fused velocity
self.x += self.vx * dt
self.y += self.vy * dt
\end{lstlisting}
\end{filebox}

\textbf{Математика предсказания:}

Команда колёс задаёт вперёд\-направленную скорость $v_{\text{cmd}}$ в теле робота.
Поворот на угол курса $\theta$ (из IMU) переводит её в мировую систему координат:

\begin{equation}
  \begin{pmatrix} v_x^{\text{wheel}} \\ v_y^{\text{wheel}} \end{pmatrix}
  = v_{\text{cmd}} \begin{pmatrix} \cos\theta \\ \sin\theta \end{pmatrix}
  \label{eq:wheel_world}
\end{equation}

Смешивание с текущим состоянием ($\alpha = 0.5$):

\begin{equation}
  \vect{x}^- = (1-\alpha)\,\vect{x} + \alpha\,\vect{x}^{\text{wheel}}
  \label{eq:blend}
\end{equation}

Рост ковариации (процессный шум $q_{\text{wheel}} = 0.01$):

\begin{equation}
  p_{00}^- = p_{00} + q_{\text{wheel}} \cdot \Delta t, \quad
  p_{11}^- = p_{11} + q_{\text{wheel}} \cdot \Delta t
  \label{eq:cov_grow}
\end{equation}

Интегрирование позиции (правило Эйлера):

\begin{equation}
  \boxed{x \mathrel{+}= v_x \cdot \Delta t, \quad y \mathrel{+}= v_y \cdot \Delta t}
  \label{eq:pos_integrate}
\end{equation}

\section{Шаг коррекции (Update) по акселерометру}
\label{sec:ekf_update}

\begin{filebox}[title={\filepath{velocity\_ekf.py}, строки 122--173 --- update\_accel()}]
\begin{lstlisting}[style=pythonstyle, numbers=left, firstnumber=135]
# Integrate accel to get velocity measurement
self._accel_vx += la_wx * dt
self._accel_vy += la_wy * dt

# Decay accel velocity to prevent unbounded drift
self._accel_vx *= 0.98
self._accel_vy *= 0.98

# Innovation: y = z - H*x  (H = I)
yx = self._accel_vx - self.vx
yy = self._accel_vy - self.vy

# S = P + R (innovation covariance)
r = self.r_accel + self.q_accel * dt
s00 = self._p00 + r
s11 = self._p11 + r

# Kalman gain: K = P * S^-1
k00 = self._p00 / s00
k11 = self._p11 / s11

# State update
self.vx += k00 * yx
self.vy += k11 * yy

# Covariance update: P = (I - K*H) * P
self._p00 *= (1 - k00)
self._p11 *= (1 - k11)
self._p01 *= (1 - 0.5 * (k00 + k11))
self._p10 *= (1 - 0.5 * (k00 + k11))
\end{lstlisting}
\end{filebox}

\textbf{Математика коррекции:}

Однократное интегрирование ускорения (измерение скорости):

\begin{equation}
  \vect{v}^{\text{accel}}_{k} = \vect{v}^{\text{accel}}_{k-1}
  + \vect{a}^{\text{world}} \cdot \Delta t
  \label{eq:accel_integrate}
\end{equation}

Экспоненциальное затухание для борьбы с дрейфом:

\begin{equation}
  \vect{v}^{\text{accel}} \leftarrow 0.98 \cdot \vect{v}^{\text{accel}}
  \label{eq:decay}
\end{equation}

Инновация (невязка):

\begin{equation}
  \vect{y} = \vect{z} - \mat{H}\,\vect{x}^-, \quad \mat{H} = \mat{I}
  \quad\Rightarrow\quad
  \vect{y} = \vect{v}^{\text{accel}} - \vect{v}^{\text{state}}
  \label{eq:innovation}
\end{equation}

Ковариация инновации:

\begin{equation}
  \mat{S} = \mat{P}^- + \mat{R}, \quad
  R = r_{\text{accel}} + q_{\text{accel}} \cdot \Delta t
  \label{eq:S_matrix}
\end{equation}

\textbf{Усиление Калмана} (для диагональной $\mat{S}$, т.е. поэлементно):

\begin{equation}
  \boxed{K_{ii} = \frac{p_{ii}^-}{S_{ii}}, \quad i \in \{0,1\}}
  \label{eq:kalman_gain}
\end{equation}

\textbf{Обновление состояния:}

\begin{equation}
  \vect{x} = \vect{x}^- + \mat{K}\,\vect{y}
  \label{eq:state_update}
\end{equation}

\textbf{Обновление ковариации} (формула Джозефа в диагональном виде):

\begin{equation}
  \boxed{p_{ii} \leftarrow (1 - K_{ii})\,p_{ii}}
  \label{eq:cov_update}
\end{equation}

\section{ZUPT --- обновление нулевой скорости}

\begin{filebox}[title={\filepath{velocity\_ekf.py}, строки 181--207 --- update\_zupt()}]
\begin{lstlisting}[style=pythonstyle, numbers=left, firstnumber=186]
# Kalman gain (measurement = zero, R = r_zupt)
s00 = self._p00 + r
s11 = self._p11 + r
k00 = self._p00 / s00
k11 = self._p11 / s11

# Update: z = [0, 0], innovation = -[vx, vy]
self.vx -= k00 * self.vx
self.vy -= k11 * self.vy

# Covariance update
self._p00 *= (1 - k00)
self._p11 *= (1 - k11)
\end{lstlisting}
\end{filebox}

Когда робот стоит, измерение $\vect{z} = \vect{0}$ с очень малым шумом
$r_{\text{zupt}} = 10^{-4}$. Инновация: $\vect{y} = \vect{0} - \vect{v} = -\vect{v}$.
Это приводит к быстрому обнулению скорости через стандартный шаг Калмана.

\section{Неголономное ограничение}
\label{sec:nonholomic}

\begin{filebox}[title={\filepath{velocity\_ekf.py}, строки 209--228 --- apply\_nonholonomic()}]
\begin{lstlisting}[style=pythonstyle, numbers=left, firstnumber=213]
cy, sy = cos(theta), sin(theta)
fwd = self.vx * cy + self.vy * sy    # forward component
lat = -self.vx * sy + self.vy * cy   # lateral component

lat *= self._lateral_decay           # kill 70% of lateral vel

self.vx = fwd * cy - lat * sy
self.vy = fwd * sy + lat * cy
\end{lstlisting}
\end{filebox}

Гусеничный робот \textbf{не может двигаться боком}. После шага Калмана
производится проекция скорости на направление движения:

\begin{align}
  v_{\text{fwd}} &= v_x \cos\theta + v_y \sin\theta \label{eq:fwd} \\
  v_{\text{lat}} &= -v_x \sin\theta + v_y \cos\theta \label{eq:lat}
\end{align}

Боковая составляющая подавляется: $v_{\text{lat}} \leftarrow 0.3 \cdot v_{\text{lat}}$.

Обратное преобразование в мировую СК (поворот матрицей $\mat{R}^{-1} = \mat{R}\T$):

\begin{equation}
  \begin{pmatrix} v_x \\ v_y \end{pmatrix}
  = \begin{pmatrix} \cos\theta & -\sin\theta \\ \sin\theta & \cos\theta \end{pmatrix}
  \begin{pmatrix} v_{\text{fwd}} \\ v_{\text{lat}} \end{pmatrix}
  \label{eq:rot_back}
\end{equation}

\section{Модуль скорости}

\begin{filebox}[title={\filepath{velocity\_ekf.py}, строки 243--245}]
\begin{lstlisting}[style=pythonstyle, numbers=left, firstnumber=243]
@property
def velocity_magnitude(self) -> float:
    return sqrt(self.vx * self.vx + self.vy * self.vy)
\end{lstlisting}
\end{filebox}

\begin{equation}
  |\vect{v}| = \sqrt{v_x^2 + v_y^2}
  \label{eq:vmag}
\end{equation}

\section{Теоретическое обоснование: связь с ТАУ}

\subsection{Модель в пространстве состояний}

Velocity EKF реализует дискретную модель в пространстве состояний
(см.\ гл.~\ref{ch:theory}, формула~\eqref{eq:state_space}):

\begin{equation}
  \underbrace{\begin{pmatrix} v_x \\ v_y \end{pmatrix}_{k}}_{\vect{x}_k}
  = \underbrace{\begin{pmatrix} 1-\alpha & 0 \\ 0 & 1-\alpha \end{pmatrix}}_{\mat{F}}
  \underbrace{\begin{pmatrix} v_x \\ v_y \end{pmatrix}_{k-1}}_{\vect{x}_{k-1}}
  + \underbrace{\begin{pmatrix} \alpha\cos\theta & 0 \\ 0 & \alpha\sin\theta \end{pmatrix}}_{\mat{B}}
  \underbrace{\begin{pmatrix} v_{\text{cmd}} \\ v_{\text{cmd}} \end{pmatrix}}_{\vect{u}}
  \label{eq:vekf_ss}
\end{equation}

\subsection{Устойчивость фильтра}

Собственные числа матрицы перехода $\mat{F}$:
$\lambda_{1,2} = 1 - \alpha = 0.5$.
Поскольку $|\lambda_i| < 1$ для дискретной системы,
фильтр \textbf{асимптотически устойчив} (ошибка сходится к нулю).

Это дискретный аналог критерия устойчивости~\eqref{eq:stability_eigenvalues}:
для дискретных систем все собственные числа должны лежать
внутри единичного круга $|\lambda_i| < 1$.

\subsection{Оптимальность усиления Калмана}

Усиление Калмана $K_{ii} = p_{ii}/(p_{ii} + R)$
минимизирует апостериорную дисперсию оценки:

\begin{equation}
  \frac{\partial}{\partial K}\,E\!\left[(x - \hat{x})^2\right] = 0
  \quad\Rightarrow\quad K^* = \frac{P^-}{P^- + R}
  \label{eq:optimal_K}
\end{equation}

При $R \to 0$ (идеальный датчик): $K \to 1$, фильтр доверяет измерению.
При $R \to \infty$ (шумный датчик): $K \to 0$, фильтр доверяет модели.

\section{Сводная таблица параметров}

\begin{center}
\begin{tabular}{llll}
\toprule
\textbf{Параметр} & \textbf{Символ} & \textbf{Значение} & \textbf{Физический смысл}\\
\midrule
\texttt{q\_wheel} & $q_w$ & $0.01\,(m/s)^2$ & Шум процесса (колёса)\\
\texttt{q\_accel} & $q_a$ & $0.05\,(m/s)^2$ & Скорость дрейфа акселерометра\\
\texttt{r\_accel} & $r_a$ & $0.1\,(m/s)^2$ & Шум измерения акселерометра\\
\texttt{r\_zupt}  & $r_z$ & $10^{-4}\,(m/s)^2$ & Доверие к нулевой скорости\\
\texttt{lateral\_decay} & --- & $0.3$ & Подавление бокового движения\\
\bottomrule
\end{tabular}
\end{center}

% ============================================================
%  Глава 2 — IMU EKF: фильтр Калмана для курсового угла
%  Файл: pi_nodes/filters/ekf_imu.py
% ============================================================
\chapter{IMU EKF: одномерный фильтр Калмана для угла курса}
\label{ch:imu_ekf}

\begin{filebox}[title={Файл исходного кода}]
\filepath{pi\_nodes/filters/ekf\_imu.py} \quad (171 строка, Python)
\end{filebox}

\section{Постановка задачи}

Наземный робот движется по плоской поверхности. Нас интересуют три угла Эйлера:
\begin{itemize}
  \item \textbf{Курс (yaw)} $\psi$ --- угол поворота вокруг вертикальной оси;
    главный навигационный параметр.
  \item \textbf{Крен (roll)} $\phi$ и \textbf{тангаж (pitch)} $\theta$ ---
    информационные, для контроля наклона на склоне.
\end{itemize}

Подход:
\begin{itemize}
  \item \textbf{Курс}: 1D-фильтр Калмана по интегрированию гироскопа $g_z$
    с отслеживанием дрейфа нуля.
  \item \textbf{Крен/тангаж}: прямой расчёт через $\mathrm{atan2}$ из акселерометра,
    без фильтрации (для наземного робота достаточно).
\end{itemize}

\section{Вектор состояния}

\begin{equation}
  \vect{x}_{\text{yaw}} = \begin{pmatrix} \psi \\ b_z \end{pmatrix}
  \label{eq:yaw_state}
\end{equation}

где $\psi$ --- угол курса (рад), $b_z$ --- дрейф нуля гироскопа по оси $Z$ (рад/с).

Матрица ковариации $2\times2$ (без numpy, три скаляра):

\begin{equation}
  \mat{P} = \begin{pmatrix} P_\psi & P_{\times} \\ P_{\times} & P_b \end{pmatrix}
  \label{eq:yaw_cov}
\end{equation}

\section{Шаг предсказания: интегрирование гироскопа}
\label{sec:imu_predict}

\begin{filebox}[title={\filepath{ekf\_imu.py}, строки 85--122 --- predict()}]
\begin{lstlisting}[style=pythonstyle, numbers=left, firstnumber=97]
# Bias-corrected gz
wz = gz - self._gz_bias

# ZUPT: if angular rate is tiny, skip integration
if abs(wz) < self._zupt_threshold:   # 0.01 rad/s
    self._P_yaw  += self.q_angle * dt * 0.1
    self._P_bias += self.q_bias * dt
    return

# Yaw prediction: simple integration
self._yaw += wz * dt

# Normalize to [-pi, pi]
self._yaw = (self._yaw + pi) % (2 * pi) - pi

# Covariance prediction: F = [[1, -dt], [0, 1]]
# P = F @ P @ F.T + Q  (manual 2x2)
p11 = self._P_yaw
p12 = self._P_cross
p22 = self._P_bias

self._P_yaw   = p11 + (-dt)*p12 + (-dt)*(p12 + (-dt)*p22) \
                + self.q_angle * dt
self._P_cross = p12 + (-dt) * p22
self._P_bias  = p22 + self.q_bias * dt
\end{lstlisting}
\end{filebox}

\textbf{Математика предсказания:}

Скорректированная угловая скорость:

\begin{equation}
  \omega_z = g_z - b_z
  \label{eq:wz_corrected}
\end{equation}

ZUPT (Zero-velocity Update): при $|\omega_z| < 0.01\,\text{рад/с}$ интегрирование
пропускается, ковариация растёт медленнее (коэффициент $0.1$).

Интегрирование угла курса:

\begin{equation}
  \psi_{k} = \psi_{k-1} + \omega_z \cdot \Delta t
  \label{eq:yaw_integrate}
\end{equation}

Нормализация в $[-\pi, \pi]$:

\begin{equation}
  \psi \leftarrow (\psi + \pi) \bmod 2\pi - \pi
  \label{eq:yaw_normalize}
\end{equation}

Матрица перехода системы:

\begin{equation}
  \mat{F} = \begin{pmatrix} 1 & -\Delta t \\ 0 & 1 \end{pmatrix}
  \label{eq:F_matrix}
\end{equation}

Уравнение предсказания ковариации:

\begin{equation}
  \boxed{\mat{P}^- = \mat{F}\,\mat{P}\,\mat{F}\T + \mat{Q}}
  \label{eq:P_predict}
\end{equation}

где матрица процессного шума:

\begin{equation}
  \mat{Q} = \begin{pmatrix} q_{\text{angle}} \cdot \Delta t & 0 \\ 0 & q_{\text{bias}} \cdot \Delta t \end{pmatrix}
  \label{eq:Q_matrix}
\end{equation}

В развёрнутом виде (ручное умножение $2\times2$):

\begin{align}
  P_\psi^- &= P_\psi - \Delta t \cdot P_{\times}
              - \Delta t (P_{\times} - \Delta t \cdot P_b)
              + q_{\text{angle}} \cdot \Delta t \label{eq:Pyaw_pred}\\
  P_{\times}^- &= P_{\times} - \Delta t \cdot P_b \label{eq:Pcross_pred}\\
  P_b^- &= P_b + q_{\text{bias}} \cdot \Delta t \label{eq:Pbias_pred}
\end{align}

\section{Шаг коррекции: крен и тангаж из акселерометра}

\begin{filebox}[title={\filepath{ekf\_imu.py}, строки 124--147 --- update()}]
\begin{lstlisting}[style=pythonstyle, numbers=left, firstnumber=132]
a_mag = sqrt(ax*ax + ay*ay + az*az)
if abs(a_mag - self.g) > self.accel_gate * self.g:
    return   # robot is accelerating -- skip gravity update

# Raw roll/pitch from gravity vector
raw_roll  = atan2(ay, az)
raw_pitch = atan2(-ax, sqrt(ay*ay + az*az))

# Subtract home position -> angles relative to startup
self._roll  = raw_roll  - self._home_roll
self._pitch = raw_pitch - self._home_pitch

# Normalize to [-pi, pi]
self._roll  = (self._roll  + pi) % (2*pi) - pi
self._pitch = (self._pitch + pi) % (2*pi) - pi
\end{lstlisting}
\end{filebox}

\textbf{Математика определения крена и тангажа:}

Когда робот не ускоряется, вектор акселерометра совпадает с вектором гравитации
$\vect{g} = (0, 0, -g)$ в мировой СК:

\begin{equation}
  |\vect{a}| = \sqrt{a_x^2 + a_y^2 + a_z^2} \approx g = 9.81\,\text{м/с}^2
  \label{eq:accel_gate}
\end{equation}

Условие пропуска обновления: $\bigl||\vect{a}| - g\bigr| > 0.3g$.

Крен (вращение вокруг оси $X$):

\begin{equation}
  \boxed{\phi_{\text{raw}} = \mathrm{atan2}(a_y,\, a_z)}
  \label{eq:roll_atan2}
\end{equation}

Тангаж (вращение вокруг оси $Y$):

\begin{equation}
  \boxed{\theta_{\text{raw}} = \mathrm{atan2}\!\left(-a_x,\, \sqrt{a_y^2 + a_z^2}\right)}
  \label{eq:pitch_atan2}
\end{equation}

Относительный угол (от положения при старте):

\begin{equation}
  \phi = \phi_{\text{raw}} - \phi_{\text{home}}, \quad
  \theta = \theta_{\text{raw}} - \theta_{\text{home}}
  \label{eq:rel_angles}
\end{equation}

\begin{notebox}
Акселерометр \textbf{не} корректирует угол курса $\psi$: для этого нужен
магнитометр. Поэтому угол курса накапливает ошибку интегрирования ---
именно поэтому используется фильтр Калмана с отслеживанием дрейфа $b_z$.
\end{notebox}

\section{Теоретическое обоснование: линеаризация и устойчивость}

\subsection{Нелинейная модель и якобиан}

Нелинейная модель гироскопа с дрейфом нуля:

\begin{equation}
  \vect{f}(\vect{x}, u) = \begin{pmatrix} \psi + (g_z - b_z)\Delta t \\ b_z \end{pmatrix}
  \label{eq:imu_nonlinear}
\end{equation}

Якобиан этой функции по вектору состояния $\vect{x} = (\psi, b_z)\T$
(метод линеаризации, гл.~\ref{ch:theory}, раздел~\ref{sec:jacobian}):

\begin{equation}
  \mat{F} = \frac{\partial \vect{f}}{\partial \vect{x}} =
  \begin{pmatrix}
    \partial f_1/\partial\psi & \partial f_1/\partial b_z \\
    \partial f_2/\partial\psi & \partial f_2/\partial b_z
  \end{pmatrix}
  = \begin{pmatrix} 1 & -\Delta t \\ 0 & 1 \end{pmatrix}
  \label{eq:imu_jacobian}
\end{equation}

\subsection{Устойчивость фильтра}

Собственные числа $\mat{F}$: $\lambda_{1,2} = 1$ (двукратный корень).
Без коррекции система \textbf{нейтрально устойчива} --- ошибка не растёт,
но и не убывает. Именно шаг коррекции (Update) через усиление Калмана
обеспечивает асимптотическую устойчивость, уменьшая ковариацию $\mat{P}$.

\subsection{Наблюдаемость}

Матрица наблюдаемости для $\mat{H} = (1, 0)$
(прямое измерение $\psi$ через внешний датчик):

\begin{equation}
  \mat{S}_o = \begin{pmatrix} \mat{H} \\ \mat{H}\mat{F} \end{pmatrix}
  = \begin{pmatrix} 1 & 0 \\ 1 & -\Delta t \end{pmatrix}, \quad
  \mathrm{rank}(\mat{S}_o) = 2
  \label{eq:imu_observability}
\end{equation}

Ранг равен размерности состояния $\Rightarrow$ система
\textbf{полностью наблюдаема}: дрейф $b_z$ может быть оценён
из измерений угла $\psi$.

\subsection{Связь с передаточной функцией}

В операторной форме модель гироскопа:
\begin{equation}
  \psi(p) = \frac{1}{p}\bigl(g_z(p) - b_z(p)\bigr)
  \label{eq:gyro_tf}
\end{equation}

Это интегрирующее звено с передаточной функцией $W(p) = 1/p$,
что объясняет накопление ошибки (дрейф) при отсутствии коррекции.

\section{Сводная таблица параметров}

\begin{center}
\begin{tabular}{llll}
\toprule
\textbf{Параметр} & \textbf{Символ} & \textbf{Значение} & \textbf{Смысл}\\
\midrule
\texttt{q\_angle} & $q_\psi$ & $0.001$ & Шум процесса по курсу\\
\texttt{q\_bias}  & $q_b$ & $0.0001$ & Шум дрейфа нуля гиро\\
\texttt{r\_accel} & $r_a$ & $0.5$ & Шум измерения акселерометра\\
\texttt{accel\_gate} & --- & $0.3$ & Порог ворот акселерометра ($30\%$)\\
\texttt{zupt\_threshold} & --- & $0.01\,\text{рад/с}$ & Порог ZUPT (${\approx}0.6\,\text{°/с}$)\\
\bottomrule
\end{tabular}
\end{center}

% ============================================================
%  Глава 3 — AccelPositionEstimator: оценка позиции
%  Файл: pi_nodes/filters/accel_position.py
% ============================================================
\chapter{Оценка позиции по акселерометру с слиянием датчиков}
\label{ch:accel_position}

\begin{filebox}[title={Файл исходного кода}]
\filepath{pi\_nodes/filters/accel\_position.py} \quad (420 строк, Python)
\end{filebox}

\section{Архитектура системы слияния датчиков}

Модуль объединяет три источника информации:

\begin{enumerate}
  \item \textbf{Акселерометр} --- мгновенное линейное ускорение (50 Гц)
  \item \textbf{Колёсная одометрия} --- команды моторов (20 Гц)
  \item \textbf{Orientation from IMU} --- углы Эйлера для поворота фреймов
\end{enumerate}

\section{Константы физики}

\begin{filebox}[title={\filepath{accel\_position.py}, строки 44--69}]
\begin{lstlisting}[style=pythonstyle, numbers=left, firstnumber=44]
EARTH_OMEGA = 7.2921e-5   # rad/s -- Earth angular velocity
DEFAULT_LATITUDE = 55.75  # Moscow

ZUPT_GYRO_ENTER  = 0.015  # rad/s
ZUPT_ACCEL_ENTER = 0.10   # m/s^2
ZUPT_GYRO_EXIT   = 0.04   # rad/s
ZUPT_ACCEL_EXIT  = 0.25   # m/s^2
ZUPT_ENTER_COUNT = 5      # 100ms @ 50Hz
ZUPT_EXIT_COUNT  = 3      # 60ms  @ 50Hz

LATERAL_DECAY = 0.5       # kill 50% lateral velocity per step
\end{lstlisting}
\end{filebox}

\textbf{Угловая скорость Земли} (для компенсации Кориолиса):

\begin{equation}
  \Omega = 7.2921 \times 10^{-5}\,\text{рад/с}
  \label{eq:earth_omega}
\end{equation}

\textbf{Проекция на широту} (Москва, $\varphi = 55.75^\circ$):

\begin{equation}
  \omega_z = \Omega \sin\varphi
  \label{eq:omega_z}
\end{equation}

\section{Шаг 1: Вычитание гравитации}

\begin{filebox}[title={\filepath{accel\_position.py}, строки 211--225 --- update\_imu()}]
\begin{lstlisting}[style=pythonstyle, numbers=left, firstnumber=211]
# Remove gravity -- DIRECT SUBTRACTION
la_bx = ax - self._g_body[0]
la_by = ay - self._g_body[1]
la_bz = az - self._g_body[2]

# Delta orientation correction (slopes only)
d_roll  = roll
d_pitch = pitch
if abs(d_roll) > 0.01 or abs(d_pitch) > 0.01:
    g = self._g_mag
    delta_gx = -g * sin(d_pitch)
    delta_gy  =  g * sin(d_roll)
    la_bx -= delta_gx
    la_by -= delta_gy
\end{lstlisting}
\end{filebox}

\textbf{Прямое вычитание} вектора гравитации, откалиброванного при старте:

\begin{equation}
  \vect{a}^{\text{lin}}_{\text{body}} = \vect{a}_{\text{meas}} - \vect{g}_{\text{body}}
  \label{eq:gravity_subtract}
\end{equation}

\begin{notebox}
\textbf{Почему прямое вычитание, а не ротационный метод?}\\
IMU физически наклонён на шасси. В покое он показывает, например,
$a_x = 0.707\,\text{м/с}^2$ --- чистая гравитация через наклон.
При вычитании: $0.707 - 0.707 = 0$ --- нет ложного движения.
Ротационный метод предполагает $\text{крен}=\text{тангаж}=0$, давая $g_{\text{body}} = (0, 0, 9.81)$,
и оставляет $0.707\,\text{м/с}^2$ как «движение».
\end{notebox}

Поправка на наклон (при нахождении на уклоне $> 0.6^\circ$):

\begin{equation}
  \Delta g_x = -g \sin\theta_{\text{pitch}}, \quad
  \Delta g_y = g \sin\phi_{\text{roll}}
  \label{eq:slope_correction}
\end{equation}

\section{Шаг 2: Поворот из тела в мировую СК}

\begin{filebox}[title={\filepath{accel\_position.py}, строки 227--230}]
\begin{lstlisting}[style=pythonstyle, numbers=left, firstnumber=227]
# Body -> World frame (rotate by yaw)
cy, sy = cos(yaw), sin(yaw)
la_wx = la_bx * cy - la_by * sy
la_wy = la_bx * sy + la_by * cy
\end{lstlisting}
\end{filebox}

Матрица поворота вокруг вертикальной оси ($z$-axis rotation):

\begin{equation}
  \boxed{
  \begin{pmatrix} a^{\text{world}}_x \\ a^{\text{world}}_y \end{pmatrix}
  = \mat{R}(\psi)\,
  \begin{pmatrix} a^{\text{body}}_x \\ a^{\text{body}}_y \end{pmatrix},
  \quad
  \mat{R}(\psi) = \begin{pmatrix} \cos\psi & -\sin\psi \\ \sin\psi & \cos\psi \end{pmatrix}
  }
  \label{eq:body_to_world}
\end{equation}

\section{Шаг 3: Компенсация эффекта Кориолиса}

\begin{filebox}[title={\filepath{accel\_position.py}, строки 232--234}]
\begin{lstlisting}[style=pythonstyle, numbers=left, firstnumber=232]
# Earth rotation compensation (Coriolis)
la_wx -= 2.0 * self._omega_z * self.vy
la_wy += 2.0 * self._omega_z * self.vx
\end{lstlisting}
\end{filebox}

Сила Кориолиса на единицу массы:

\begin{equation}
  \vect{a}_{\text{Cor}} = -2\,\vect{\Omega} \times \vect{v}
  \label{eq:coriolis}
\end{equation}

В 2D (только вертикальная проекция $\omega_z = \Omega\sin\varphi$):

\begin{equation}
  \begin{cases}
    a_{Cx} = -2\,\omega_z\,v_y \\
    a_{Cy} = +2\,\omega_z\,v_x
  \end{cases}
  \label{eq:coriolis_2d}
\end{equation}

Вычитание ускорения Кориолиса из измеренного:

\begin{equation}
  a^{\text{lin}}_x \leftarrow a^{\text{lin}}_x - 2\omega_z v_y, \quad
  a^{\text{lin}}_y \leftarrow a^{\text{lin}}_y + 2\omega_z v_x
  \label{eq:coriolis_remove}
\end{equation}

\section{Шаг 4: Адаптивный медианный фильтр}

\begin{filebox}[title={\filepath{accel\_position.py}, строки 236--250}]
\begin{lstlisting}[style=pythonstyle, numbers=left, firstnumber=236]
# Median filter (reject vibration spikes)
self._buf_ax.append(la_wx)
self._buf_ay.append(la_wy)

if len(self._buf_ax) >= 3:
    la_wx = _median_of_3(
        self._buf_ax[-1], self._buf_ax[-2], self._buf_ax[-3])
    la_wy = _median_of_3(
        self._buf_ay[-1], self._buf_ay[-2], self._buf_ay[-3])

# Noise threshold -- zero out tiny accelerations
if abs(la_wx) < ACCEL_NOISE_THRESHOLD:   # 0.15 m/s^2
    la_wx = 0.0
\end{lstlisting}
\end{filebox}

Медиана трёх последних значений (без сортировки --- одно ветвление):

\begin{equation}
  \hat{a} = \mathrm{median}(a_{k},\, a_{k-1},\, a_{k-2})
  \label{eq:median3}
\end{equation}

Пороговое обнуление шума: $|a| < 0.15\,\text{м/с}^2 \Rightarrow a = 0$.

\section{Шаг 5: ZUPT с гистерезисом}

\begin{filebox}[title={\filepath{accel\_position.py}, строки 252--305}]
\begin{lstlisting}[style=pythonstyle, numbers=left, firstnumber=252]
gyro_mag  = sqrt(gx*gx + gy*gy + gz*gz)
accel_mag = sqrt(la_wx*la_wx + la_wy*la_wy)

if self._stationary:
    # Need STRONG evidence of motion to exit
    if (self._cmd_moving and
        (gyro_mag > ZUPT_GYRO_EXIT or       # 0.04 rad/s
         accel_mag > ZUPT_ACCEL_EXIT)):      # 0.25 m/s^2
        self._exit_count += 1
    if self._exit_count >= ZUPT_EXIT_COUNT:  # 3 samples
        self._stationary = False
else:
    # Easy to re-enter stationary
    if (not self._cmd_moving or
        (gyro_mag < ZUPT_GYRO_ENTER and     # 0.015 rad/s
         accel_mag < ZUPT_ACCEL_ENTER)):    # 0.10 m/s^2
        self._enter_count += 1
    if self._enter_count >= ZUPT_ENTER_COUNT: # 5 samples
        self._stationary = True
\end{lstlisting}
\end{filebox}

\textbf{Двухпороговый автомат состояний (гистерезис):}

\begin{center}
\begin{tikzpicture}[
  state/.style={rectangle, rounded corners, draw=blue!60, fill=blue!10,
                minimum width=3cm, minimum height=1.2cm, font=\bfseries},
  arrow/.style={-{Stealth[length=8pt]}, thick}
]
  \node[state] (S) {STATIONARY};
  \node[state, right=5cm of S] (M) {MOVING};

  \draw[arrow, bend left=20] (S) to node[above, font=\small, align=center]{
    $|\omega|>0.04$ \textbf{или} $|a|>0.25$ м/с$^2$\\
    \textit{за 3 последовательных шага}
  } (M);

  \draw[arrow, bend left=20] (M) to node[below, font=\small, align=center]{
    $|\omega|<0.015$ \textbf{и} $|a|<0.10$ м/с$^2$\\
    \textit{за 5 последовательных шагов}
  } (S);
\end{tikzpicture}
\end{center}

Высокие пороги \textit{выхода} ($> 0.04, > 0.25$) и низкие пороги \textit{входа}
($< 0.015, < 0.10$) предотвращают ложные переходы на границе движения/покоя.

\section{Шаг 6: Трапецеидальное интегрирование (резервный режим)}

\begin{filebox}[title={\filepath{accel\_position.py}, строки 315--334 --- fallback}]
\begin{lstlisting}[style=pythonstyle, numbers=left, firstnumber=315]
# Trapezoidal integration (more accurate than Euler)
self.vx += 0.5 * (la_wx + self._prev_la_x) * dt
self.vy += 0.5 * (la_wy + self._prev_la_y) * dt

# Non-holonomic constraint
fwd = self.vx * cy + self.vy * sy
lat = -self.vx * sy + self.vy * cy
lat *= LATERAL_DECAY
self.vx = fwd * cy - lat * sy
self.vy = fwd * sy + lat * cy

# Position integration
self._accel_x += self.vx * dt
self._accel_y += self.vy * dt
\end{lstlisting}
\end{filebox}

\textbf{Метод трапеций} (порядок точности $O(\Delta t^2)$, лучше Эйлера $O(\Delta t)$):

\begin{equation}
  v_{k} = v_{k-1} + \frac{\Delta t}{2}\bigl(a_{k} + a_{k-1}\bigr)
  \label{eq:trapezoidal}
\end{equation}

\section{Шаг 7: Комплементарный фильтр (резервный режим)}

\begin{filebox}[title={\filepath{accel\_position.py}, строки 376--386 --- blend()}]
\begin{lstlisting}[style=pythonstyle, numbers=left, firstnumber=378]
# Complementary filter: alpha*wheel + (1-alpha)*accel
self._alpha = 0.95 * self._alpha + 0.05 * self._alpha_moving
a = self._alpha
self.x = a * self._wheel_x + (1.0 - a) * self._accel_x
self.y = a * self._wheel_y + (1.0 - a) * self._accel_y

# Slow drift correction of accel toward wheel
self._accel_x = 0.995 * self._accel_x + 0.005 * self._wheel_x
self._accel_y = 0.995 * self._accel_y + 0.005 * self._wheel_y
\end{lstlisting}
\end{filebox}

Комплементарный фильтр с адаптивным $\alpha$:

\begin{equation}
  \vect{p} = \alpha\,\vect{p}_{\text{wheel}} + (1-\alpha)\,\vect{p}_{\text{accel}}
  \label{eq:complementary}
\end{equation}

где $\alpha \in [0.95, 1.0]$ (при движении растёт вес одометрии).
Медленная коррекция дрейфа акселерометра:
$\vect{p}_{\text{accel}} \leftarrow 0.995\,\vect{p}_{\text{accel}} + 0.005\,\vect{p}_{\text{wheel}}$.

\section{Теоретическое обоснование: слияние датчиков}

\subsection{Комплементарный фильтр как частный случай фильтра Калмана}

Комплементарный фильтр~\eqref{eq:complementary} можно рассматривать
как \textbf{стационарный фильтр Калмана} с фиксированным усилением:

\begin{equation}
  \hat{x} = \alpha\,x_{\text{odom}} + (1-\alpha)\,x_{\text{accel}}
  = x_{\text{accel}} + \alpha\,(x_{\text{odom}} - x_{\text{accel}})
  \label{eq:comp_as_kalman}
\end{equation}

Сравнивая с~\eqref{eq:kf_update_x}: $\hat{x} = \hat{x}^- + K\cdot y$,
видим что $\alpha$ играет роль \textbf{усиления Калмана} $K$.

\subsection{Частотная интерпретация}

В частотной области комплементарный фильтр выделяет:
\begin{itemize}
  \item \textbf{Низкие частоты} из одометрии (стабильная, но дрейфующая)
  \item \textbf{Высокие частоты} из акселерометра (шумный, но без дрейфа)
\end{itemize}

\begin{equation}
  X(p) = \underbrace{\frac{\alpha}{1 + \alpha/p}}_{\text{ФНЧ}} X_{\text{odom}}(p)
       + \underbrace{\frac{p/\alpha}{1 + p/\alpha}}_{\text{ФВЧ}} X_{\text{accel}}(p)
  \label{eq:comp_frequency}
\end{equation}

Частота среза: $f_c = \alpha/(2\pi\Delta t) \approx 0.95/(2\pi \cdot 0.02) \approx 7.6\,\text{Гц}$.

\subsection{Трапецеидальное vs.\ Эйлерово интегрирование}

\begin{center}
\begin{tabular}{lll}
\toprule
\textbf{Метод} & \textbf{Формула} & \textbf{Погрешность}\\
\midrule
Эйлер вперёд & $v_k = v_{k-1} + a_k \Delta t$ & $O(\Delta t)$\\
Трапеция & $v_k = v_{k-1} + \frac{a_k + a_{k-1}}{2}\Delta t$ & $O(\Delta t^2)$\\
\bottomrule
\end{tabular}
\end{center}

При частоте $50\,\text{Гц}$ ($\Delta t = 0.02\,\text{с}$) разница в погрешности
составляет $\Delta t/2 = 0.01$, т.е.\ трапеция точнее Эйлера в ${\approx}100$ раз
для гладких сигналов.

\subsection{Дискретная модель пространства состояний}

Полная система AccelPosition в форме пространства состояний
(гл.~\ref{ch:theory}, формула~\eqref{eq:discrete_ss}):

\begin{equation}
  \underbrace{\begin{pmatrix} x \\ y \\ v_x \\ v_y \end{pmatrix}_{k}}_{\vect{x}_k}
  = \underbrace{\begin{pmatrix}
    1 & 0 & \Delta t & 0 \\
    0 & 1 & 0 & \Delta t \\
    0 & 0 & 1 & 0 \\
    0 & 0 & 0 & 1
  \end{pmatrix}}_{\mat{F}}
  \vect{x}_{k-1}
  + \underbrace{\begin{pmatrix}
    \Delta t^2/2 & 0 \\
    0 & \Delta t^2/2 \\
    \Delta t & 0 \\
    0 & \Delta t
  \end{pmatrix}}_{\mat{B}}
  \begin{pmatrix} a_x \\ a_y \end{pmatrix}
  \label{eq:accel_ss}
\end{equation}

\section{Сводная схема конвейера обработки}

\begin{center}
\begin{tikzpicture}[
  node distance=0.6cm,
  block/.style={rectangle, draw=blue!50, fill=blue!8, rounded corners,
                minimum width=4.5cm, minimum height=0.8cm,
                font=\small\bfseries, align=center},
  arrow/.style={-{Stealth}, thick, blue!60}
]
  \node[block] (raw) {Акселерометр \\ $\vect{a}_{\text{meas}}$};
  \node[block, below=of raw] (grav) {Вычитание гравитации \\ $\vect{a}^{\text{lin}} = \vect{a} - \vect{g}_{\text{body}}$};
  \node[block, below=of grav] (rot)  {Поворот СК \\ $\vect{a}^{\text{world}} = \mat{R}(\psi)\vect{a}^{\text{lin}}$};
  \node[block, below=of rot]  (cor)  {Компенсация Кориолиса};
  \node[block, below=of cor]  (med)  {Медианный фильтр + порог шума};
  \node[block, below=of med]  (zupt) {ZUPT с гистерезисом};
  \node[block, below=of zupt] (ekf)  {Velocity EKF \\ (слияние с одометрией)};
  \node[block, below=of ekf]  (nh)   {Неголономное ограничение};
  \node[block, below=of nh]   (pos)  {Позиция $(x, y)$};

  \draw[arrow] (raw)  -- (grav);
  \draw[arrow] (grav) -- (rot);
  \draw[arrow] (rot)  -- (cor);
  \draw[arrow] (cor)  -- (med);
  \draw[arrow] (med)  -- (zupt);
  \draw[arrow] (zupt) -- (ekf);
  \draw[arrow] (ekf)  -- (nh);
  \draw[arrow] (nh)   -- (pos);
\end{tikzpicture}
\end{center}

% ============================================================
%  Глава 4 — Pure Pursuit: алгоритм следования по пути
%  Файл: pi_nodes/nodes/path_recorder_node.py
% ============================================================
\chapter{Pure Pursuit: алгоритм возврата домой}
\label{ch:pure_pursuit}

\begin{filebox}[title={Файл исходного кода}]
\filepath{pi\_nodes/nodes/path\_recorder\_node.py} \quad (342 строки, Python)
\end{filebox}

\section{Постановка задачи}

Модуль записывает путь робота в виде набора опорных точек и затем проводит
робота обратно по тому же пути в обратном порядке (возврат домой).

\begin{itemize}
  \item \textbf{Запись}: каждые $3\,\text{см}$ или $4.5^\circ$ поворота.
  \item \textbf{Воспроизведение}: алгоритм Pure Pursuit с финальным фазой
    точного наведения.
\end{itemize}

\section{Запись путевых точек}

\begin{filebox}[title={\filepath{path\_recorder\_node.py}, строки 135--148 --- \_maybe\_record\_waypoint()}]
\begin{lstlisting}[style=pythonstyle, numbers=left, firstnumber=139]
dx = self._x - self._last_record_x
dy = self._y - self._last_record_y
dist = math.sqrt(dx * dx + dy * dy)
dtheta = abs(self._angle_diff(self._theta, self._last_record_theta))

if dist >= RECORD_INTERVAL_M or dtheta >= RECORD_INTERVAL_RAD:
    self._path.append((self._x, self._y, self._theta))
    # RECORD_INTERVAL_M = 0.03 m, RECORD_INTERVAL_RAD = 0.08 rad
\end{lstlisting}
\end{filebox}

\textbf{Расстояние между точками} (Евклидова метрика):

\begin{equation}
  d = \sqrt{(x - x_{\text{last}})^2 + (y - y_{\text{last}})^2}
  \label{eq:waypoint_dist}
\end{equation}

Запись происходит при: $d \geq 0.03\,\text{м}$ ИЛИ $|\Delta\theta| \geq 0.08\,\text{рад}$.

\section{Алгоритм Pure Pursuit}
\label{sec:pure_pursuit_algo}

\begin{filebox}[title={\filepath{path\_recorder\_node.py}, строки 206--243 --- \_control\_loop()}]
\begin{lstlisting}[style=pythonstyle, numbers=left, firstnumber=206]
# Pure pursuit: simultaneous linear + angular
bearing = math.atan2(dy, dx)
angle_error = self._angle_diff(bearing, self._theta)

if in_final_approach:
    angular   = FINAL_ANGULAR_KP * angle_error    # Kp = 3.0
    angular   = max(-REPLAY_ANGULAR_SPEED, min(REPLAY_ANGULAR_SPEED, angular))
    cos_factor = max(0.0, math.cos(angle_error))
    dist_ratio = dist / FINAL_APPROACH_RADIUS      # 0..1
    linear     = FINAL_LINEAR_SPEED * cos_factor * dist_ratio
    linear     = max(REPLAY_MIN_LINEAR * 0.6 * cos_factor, linear)
    if abs(angle_error) > 0.8:
        linear = 0.0
else:
    angular    = REPLAY_ANGULAR_KP * angle_error  # Kp = 2.5
    angular    = max(-REPLAY_ANGULAR_SPEED, min(REPLAY_ANGULAR_SPEED, angular))
    cos_factor = max(0.0, math.cos(angle_error))
    dist_factor = min(1.0, dist / 0.15)
    linear     = REPLAY_LINEAR_SPEED * cos_factor * dist_factor
    linear     = max(REPLAY_MIN_LINEAR * cos_factor, linear)
    if abs(angle_error) > 1.2:
        linear = 0.0
\end{lstlisting}
\end{filebox}

\subsection{Вычисление угла на цель}

Пеленг (bearing) на текущую целевую точку $(t_x, t_y)$:

\begin{equation}
  \beta = \mathrm{atan2}(t_y - y,\; t_x - x)
  \label{eq:bearing}
\end{equation}

\textbf{Угловая ошибка} (кратчайший путь в $[-\pi, \pi]$):

\begin{equation}
  e_\theta = \mathrm{angle\_diff}(\beta, \psi) \in [-\pi, \pi]
  \label{eq:angle_error}
\end{equation}

\subsection{Пропорциональный угловой регулятор}

\begin{equation}
  \omega = K_P \cdot e_\theta, \quad
  \omega \in [-\omega_{\max}, +\omega_{\max}]
  \label{eq:angular_control}
\end{equation}

\begin{center}
\begin{tabular}{lll}
\toprule
\textbf{Режим} & $K_P$ & $\omega_{\max}$\\
\midrule
Обычный Pure Pursuit & $2.5$ & $0.8\,\text{рад/с}$\\
Финальный подход & $3.0$ & $0.8\,\text{рад/с}$\\
\bottomrule
\end{tabular}
\end{center}

\subsection{Косинусное масштабирование линейной скорости}

\textbf{Ключевая идея Pure Pursuit}: при большой угловой ошибке снижать
линейную скорость через косинус:

\begin{equation}
  \cos\text{-factor} = \max\bigl(0,\, \cos(e_\theta)\bigr)
  \label{eq:cos_factor}
\end{equation}

\begin{equation}
  v = v_{\max} \cdot \cos\text{-factor} \cdot d_{\text{factor}}
  \label{eq:linear_speed}
\end{equation}

Коэффициент дистанции:

\begin{equation}
  d_{\text{factor}} = \min\!\left(1,\, \frac{d}{0.15\,\text{м}}\right)
  \label{eq:dist_factor}
\end{equation}

При $|e_\theta| > \pi/2$: $\cos(e_\theta) \leq 0 \Rightarrow v = 0$
(только поворот на месте).

\subsection{Финальный подход}

В радиусе $0.12\,\text{м}$ от конечной точки:

\begin{equation}
  v = v_{\text{final}} \cdot \cos(e_\theta) \cdot \frac{d}{r_{\text{final}}}
  \label{eq:final_approach}
\end{equation}

Скорость пропорциональна оставшемуся расстоянию --- плавное торможение.

\begin{center}
\begin{tabular}{lll}
\toprule
\textbf{Параметр} & \textbf{Значение} & \textbf{Описание}\\
\midrule
\texttt{REPLAY\_LINEAR\_SPEED}   & $0.20\,\text{м/с}$ & Скорость воспроизведения\\
\texttt{FINAL\_LINEAR\_SPEED}    & $0.08\,\text{м/с}$ & Скорость финального подхода\\
\texttt{REPLAY\_LOOKAHEAD\_M}    & $0.10\,\text{м}$   & Дистанция look-ahead\\
\texttt{REPLAY\_HOME\_TOLERANCE} & $0.015\,\text{м}$  & Допуск финальной точки\\
\texttt{REPLAY\_GOAL\_TOLERANCE} & $0.04\,\text{м}$   & Допуск промежуточных точек\\
\bottomrule
\end{tabular}
\end{center}

\section{Нормализация угла}

\begin{filebox}[title={\filepath{path\_recorder\_node.py}, строки 312--320 --- \_angle\_diff()}]
\begin{lstlisting}[style=pythonstyle, numbers=left, firstnumber=313]
@staticmethod
def _angle_diff(target, current):
    """Shortest signed angle difference, in [-pi, pi]."""
    d = target - current
    while d > math.pi:
        d -= 2.0 * math.pi
    while d < -math.pi:
        d += 2.0 * math.pi
    return d
\end{lstlisting}
\end{filebox}

\begin{equation}
  \Delta\theta = \bigl[(\beta - \psi + \pi) \bmod 2\pi\bigr] - \pi
  \label{eq:angle_normalize}
\end{equation}

Это гарантирует, что угловая ошибка всегда в диапазоне $[-\pi, \pi]$,
а поворот всегда происходит по кратчайшему пути.

\section{Геометрия алгоритма Pure Pursuit}

\begin{center}
\begin{tikzpicture}[scale=1.4]
  % Robot
  \filldraw[blue!70] (0,0) circle (0.12);
  \draw[-{Stealth}, blue!70, thick] (0,0) -- (1.0, 0)
    node[right, font=\small]{$\psi$ (курс)};
  \node[below, font=\small, blue!70] at (0,-0.15) {Робот};

  % Target point
  \filldraw[red!70] (1.5, 1.2) circle (0.08);
  \node[right, font=\small, red!70] at (1.58, 1.2) {Цель $(t_x, t_y)$};

  % Bearing line
  \draw[-{Stealth}, red!50, dashed, thick] (0,0) -- (1.5, 1.2)
    node[midway, above left, font=\small]{$d$};

  % Angle annotation
  \draw[->] (0.5, 0) arc[start angle=0, end angle=38.7, radius=0.5];
  \node[font=\small] at (0.72, 0.2) {$e_\theta$};

  % Bearing angle label
  \draw[->] (0.3, 0) arc[start angle=0, end angle=38.7, radius=0.3];
  \node[font=\tiny, gray] at (0.48, 0.08) {$\beta$};
\end{tikzpicture}
\end{center}

$$\beta = \mathrm{atan2}(t_y - y, t_x - x), \quad e_\theta = \beta - \psi$$

\section{Теоретическое обоснование: замкнутая система управления}

\subsection{Структура контура управления}

Алгоритм Pure Pursuit реализует \textbf{замкнутый P-регулятор}
(гл.~\ref{ch:theory}, раздел~\ref{sec:p_controller}) для угловой
скорости:

\begin{equation}
  \omega = K_P \cdot e_\theta, \quad e_\theta = \beta - \psi
  \label{eq:pp_closed_loop}
\end{equation}

Передаточная функция замкнутой системы <<регулятор + робот>>
(робот как интегратор $\dot{\psi} = \omega$):

\begin{equation}
  W_{\text{зс}}(p) = \frac{K_P / p}{1 + K_P / p} = \frac{K_P}{p + K_P}
  \label{eq:pp_tf}
\end{equation}

Полюс замкнутой системы: $\lambda = -K_P = -2.5$,
что даёт экспоненциальную сходимость $e_\theta(t) = e_\theta(0)\,e^{-2.5t}$
с постоянной времени $\tau = 1/K_P = 0.4\,\text{с}$.

\subsection{Нелинейное масштабирование и устойчивость}

Косинусное масштабирование $v = v_{\max}\cos(e_\theta)$ создаёт
\textbf{нелинейную связь} между угловой ошибкой и линейной скоростью.
Это эквивалентно функции Ляпунова:

\begin{equation}
  V(e_\theta) = 1 - \cos(e_\theta) \geq 0
  \label{eq:pp_lyapunov}
\end{equation}

При $|e_\theta| < \pi/2$: $\dot{V} = \sin(e_\theta)\cdot\dot{e}_\theta < 0$
(при правильном знаке $K_P$), что гарантирует
\textbf{асимптотическую устойчивость} по Ляпунову.

\subsection{Связь с кинематической моделью}

Полная нелинейная модель робота с Pure Pursuit
(модель unicycle из гл.~\ref{ch:odometry}):

\begin{equation}
  \begin{cases}
    \dot{x} = v\cos\theta \cdot \cos(e_\theta) \cdot d_{\text{factor}} \\
    \dot{y} = v\sin\theta \cdot \cos(e_\theta) \cdot d_{\text{factor}} \\
    \dot{\theta} = K_P \cdot e_\theta
  \end{cases}
  \label{eq:pp_full_model}
\end{equation}

Это нелинейная система вида $\dot{\vect{x}} = \vect{f}(\vect{x})$
(гл.~\ref{ch:theory}, формула~\eqref{eq:nonlinear_system}),
линеаризация которой в окрестности $e_\theta = 0$ даёт
устойчивую линейную систему с полюсом $\lambda = -K_P$.

% ============================================================
%  Глава 5 — SLAM: вероятностная карта занятости
%  Файл: pi_nodes/nodes/slam_map_node.py
% ============================================================
\chapter{SLAM: вероятностная карта занятости на log-odds}
\label{ch:slam_map}

\begin{filebox}[title={Файл исходного кода}]
\filepath{pi\_nodes/nodes/slam\_map\_node.py} \quad (Python)
\end{filebox}

\section{Вероятностная модель карты}

Карта занятости представляет каждую ячейку сетки как бинарную случайную
величину: свободно ($\mathtt{F}$) или занято ($\mathtt{O}$).

\subsection{Log-odds представление}

Вместо хранения вероятности $P(\mathtt{O})$ используется \textbf{логарифм
отношения правдоподобий} (log-odds), устойчивый к крайним значениям $0$ и $1$:

\begin{equation}
  l = \log\frac{P(\mathtt{O})}{1 - P(\mathtt{O})} = \log\frac{P(\mathtt{O})}{P(\mathtt{F})}
  \label{eq:log_odds_def}
\end{equation}

Обратное преобразование:

\begin{equation}
  P(\mathtt{O}) = \frac{e^l}{1 + e^l} = \sigma(l)
  \label{eq:log_odds_inv}
\end{equation}

\textbf{Преимущества}: обновление --- простое сложение; нет умножения вероятностей.

\begin{filebox}[title={\filepath{slam\_map\_node.py}, строки 42--48 --- константы log-odds}]
\begin{lstlisting}[style=pythonstyle, numbers=left, firstnumber=42]
L_OCC  =  0.85   # log-odds increment for occupied cell
L_FREE = -0.40   # log-odds increment for free cell
L_MIN  = -2.0    # minimum log-odds (very confident free)
L_MAX  =  3.5    # maximum log-odds (very confident occupied)
L_THRESH_OCC  =  0.6   # above this -> occupied (output)
L_THRESH_FREE = -0.5   # below this -> free (output)
\end{lstlisting}
\end{filebox}

\textbf{Физический смысл значений log-odds:}

\begin{center}
\begin{tabular}{lll}
\toprule
\textbf{log-odds $l$} & \textbf{$P(\mathtt{O})$} & \textbf{Интерпретация}\\
\midrule
$-2.0$ & $11.9\%$ & Очень уверенно свободно\\
$-0.5$ & $37.8\%$ & Скорее свободно\\
$0.0$  & $50.0\%$ & Неизвестно\\
$+0.6$ & $64.6\%$ & Скорее занято\\
$+3.5$ & $97.1\%$ & Очень уверенно занято\\
\bottomrule
\end{tabular}
\end{center}

\section{Трассировка лучей (Ray Tracing)}

\begin{filebox}[title={\filepath{slam\_map\_node.py}, строки 256--285 --- \_trace\_ray()}]
\begin{lstlisting}[style=pythonstyle, numbers=left, firstnumber=256]
def _trace_ray(self, ox, oy, angle, range_m):
    cos_a = math.cos(angle)
    sin_a = math.sin(angle)

    step = CELL_SIZE_M * 0.7   # 3.5 cm step
    dist = 0.0
    hit_range = min(range_m, SENSOR_MAX_RANGE)

    # Mark cells along ray as FREE
    while dist < hit_range - OBSTACLE_THICKNESS:
        wx = ox + cos_a * dist
        wy = oy + sin_a * dist
        ci, cj = self._world_to_cell(wx, wy)
        if 0 <= ci < MAP_CELLS and 0 <= cj < MAP_CELLS:
            idx = cj * MAP_CELLS + ci
            self._grid[idx] = max(L_MIN, self._grid[idx] + L_FREE)
        dist += step

    # Mark obstacle cells as OCCUPIED
    if range_m < SENSOR_MAX_RANGE - 0.05:
        for d in range(3):
            dd = hit_range + d * CELL_SIZE_M * 0.5
            wx = ox + cos_a * dd
            wy = oy + sin_a * dd
            ci, cj = self._world_to_cell(wx, wy)
            if 0 <= ci < MAP_CELLS and 0 <= cj < MAP_CELLS:
                idx = cj * MAP_CELLS + ci
                self._grid[idx] = min(L_MAX, self._grid[idx] + L_OCC)
\end{lstlisting}
\end{filebox}

\textbf{Параметрическое уравнение луча:}

\begin{equation}
  \begin{pmatrix} w_x(t) \\ w_y(t) \end{pmatrix}
  = \begin{pmatrix} o_x \\ o_y \end{pmatrix}
  + t \begin{pmatrix} \cos\alpha \\ \sin\alpha \end{pmatrix}, \quad t \in [0, r]
  \label{eq:ray_parametric}
\end{equation}

\textbf{Байесовское обновление} (вычитание для свободных, прибавление для занятых):

\begin{equation}
  \boxed{l_{\text{new}} = \mathrm{clamp}\!\left(l + \Delta l,\; L_{\min},\; L_{\max}\right)}
  \label{eq:log_odds_update}
\end{equation}

где:
\[
  \Delta l = \begin{cases}
    L_{\text{FREE}} = -0.40 & \text{луч прошёл через ячейку}\\
    L_{\text{OCC}} = +0.85 & \text{луч достиг препятствия}
  \end{cases}
\]

\section{Ультразвуковой конус}

\begin{filebox}[title={\filepath{slam\_map\_node.py}, строки 247--254 --- сканирование конуса}]
\begin{lstlisting}[style=pythonstyle, numbers=left, firstnumber=247]
# Ultrasonic cone: trace multiple rays within the FOV
angle_start = robot_theta - SENSOR_FOV_RAD   # theta - 15 deg
angle_end   = robot_theta + SENSOR_FOV_RAD   # theta + 15 deg

angle = angle_start
while angle <= angle_end:
    self._trace_ray(robot_x, robot_y, angle, range_m)
    angle += ANGULAR_RESOLUTION              # step = 0.05 rad ~ 3 deg
\end{lstlisting}
\end{filebox}

Количество лучей на одно измерение:

\begin{equation}
  N_{\text{rays}} = \left\lfloor \frac{2 \cdot \text{SENSOR\_FOV\_RAD}}{\text{ANGULAR\_RESOLUTION}} \right\rfloor
  = \left\lfloor \frac{2 \times 0.26}{0.05} \right\rfloor = 10\,\text{ лучей}
  \label{eq:ray_count}
\end{equation}

\section{Преобразование координат}

\begin{filebox}[title={\filepath{slam\_map\_node.py}, строки 288--298 --- world $\leftrightarrow$ grid}]
\begin{lstlisting}[style=pythonstyle, numbers=left, firstnumber=288]
def _world_to_cell(self, wx, wy):
    ci = int((wx - self._origin_x) / CELL_SIZE_M)  # column
    cj = int((wy - self._origin_y) / CELL_SIZE_M)  # row
    return ci, cj

def _cell_to_world(self, ci, cj):
    wx = self._origin_x + (ci + 0.5) * CELL_SIZE_M
    wy = self._origin_y + (cj + 0.5) * CELL_SIZE_M
    return wx, wy
\end{lstlisting}
\end{filebox}

\textbf{Мировые координаты $\to$ ячейка сетки:}

\begin{equation}
  c_i = \left\lfloor \frac{w_x - o_x}{\delta} \right\rfloor, \quad
  c_j = \left\lfloor \frac{w_y - o_y}{\delta} \right\rfloor
  \label{eq:world_to_grid}
\end{equation}

\textbf{Ячейка $\to$ центр в мировых координатах:}

\begin{equation}
  w_x = o_x + (c_i + 0.5)\,\delta, \quad
  w_y = o_y + (c_j + 0.5)\,\delta
  \label{eq:grid_to_world}
\end{equation}

где $o_x = o_y = -5\,\text{м}$ (начало координат карты), $\delta = 0.05\,\text{м/ячейку}$.

\section{Кластеризация объектов с Евклидовой метрикой}

\begin{filebox}[title={\filepath{slam\_map\_node.py}, строки 170--188 --- \_detections\_cb()}]
\begin{lstlisting}[style=pythonstyle, numbers=left, firstnumber=172]
# Search for nearest existing object (same class+colour)
for obj_id, obj in self._objects.items():
    dx = obj['x'] - wx
    dy = obj['y'] - wy
    if math.sqrt(dx*dx + dy*dy) < OBJ_CLUSTER_RADIUS_M:  # 0.30 m
        matched_id = obj_id
        break

# Update position: sliding average
if matched_id:
    alpha = 0.3
    obj['x'] = round(alpha * wx + (1 - alpha) * obj['x'], 3)
    obj['y'] = round(alpha * wy + (1 - alpha) * obj['y'], 3)
\end{lstlisting}
\end{filebox}

Проверка принадлежности к кластеру:

\begin{equation}
  d_{\text{Euclid}} = \sqrt{(x_{\text{new}} - x_{\text{obj}})^2 + (y_{\text{new}} - y_{\text{obj}})^2}
  < 0.30\,\text{м}
  \label{eq:cluster_dist}
\end{equation}

Скользящее среднее позиции (EMA с $\alpha = 0.3$):

\begin{equation}
  \vect{p}_{\text{obj}} \leftarrow 0.3\,\vect{p}_{\text{new}} + 0.7\,\vect{p}_{\text{obj}}
  \label{eq:ema_position}
\end{equation}

\section{Параметры сетки}

\begin{center}
\begin{tabular}{lll}
\toprule
\textbf{Параметр} & \textbf{Значение} & \textbf{Описание}\\
\midrule
\texttt{MAP\_SIZE\_M}      & $10\,\text{м}$ & Размер карты ($\pm 5\,\text{м}$)\\
\texttt{CELL\_SIZE\_M}     & $0.05\,\text{м}$ & Размер ячейки (5 см)\\
\texttt{MAP\_CELLS}        & $200 \times 200$ & Размер сетки\\
\texttt{SENSOR\_FOV\_RAD}  & $0.26\,\text{рад}$ & Полу-угол конуса (${\approx}15^\circ$)\\
\texttt{SENSOR\_MAX\_RANGE}& $2.0\,\text{м}$ & Макс. дальность HC-SR04\\
\bottomrule
\end{tabular}
\end{center}

% ============================================================
%  Глава 6 — Одометрия: кинематика дифференциального привода
%  Файлы: ros_ws/.../motor_node.cpp  и  pi_nodes/nodes/motor_node.py
% ============================================================
\chapter{Одометрия: кинематика дифференциального привода}
\label{ch:odometry}

\begin{filebox}[title={Файлы исходного кода}]
\filepath{ros\_ws/src/robot\_pkg\_cpp/src/motor\_node.cpp} \quad (C++, строки 244--302)\\
\filepath{pi\_nodes/nodes/motor\_node.py} \quad (Python)
\end{filebox}

\section{Модель дифференциального привода}

Гусеничный робот управляется двумя независимыми бортами (левый/правый).
Линейная и угловая скорости центра масс:

\begin{equation}
  v_{\text{linear}} \in [-v_{\max}, v_{\max}], \quad
  \omega_{\text{angular}} \in [-\omega_{\max}, \omega_{\max}]
  \label{eq:velocity_limits}
\end{equation}

\section{Микшер дифференциального привода}

\begin{filebox}[title={\filepath{motor\_node.cpp}, строки 251--266 --- controlLoop()}]
\begin{lstlisting}[style=cppstyle, numbers=left, firstnumber=251]
// Convert m/s -> throttle in [-1, 1]
double lin_t = (max_lin_ > 0.0) ? (linear_  / max_lin_) : 0.0;
double ang_t = (max_ang_ > 0.0) ? (angular_ / max_ang_) : 0.0;

// Differential-drive mixer
double left  = lin_t + ang_t;
double right = lin_t - ang_t;

// Normalise: keep ratio, clamp to [-1, 1]
double max_val = std::max({std::abs(left), std::abs(right), 1.0});
if (max_val > 1.0) { left /= max_val; right /= max_val; }
\end{lstlisting}
\end{filebox}

\textbf{Формулы микшера} (нормированные команды $\in [-1, 1]$):

\begin{equation}
  \text{left} = \ell_{\text{lin}} + \ell_{\text{ang}}, \quad
  \text{right} = \ell_{\text{lin}} - \ell_{\text{ang}}
  \label{eq:diff_mixer}
\end{equation}

Нормировка для сохранения соотношения:

\begin{equation}
  m = \max\!\bigl(1,\, |\text{left}|,\, |\text{right}|\bigr), \quad
  \text{left} \leftarrow \frac{\text{left}}{m}, \quad
  \text{right} \leftarrow \frac{\text{right}}{m}
  \label{eq:normalize}
\end{equation}

\section{Интегрирование одометрии (кинематические уравнения)}

\begin{filebox}[title={\filepath{motor\_node.cpp}, строки 268--273}]
\begin{lstlisting}[style=cppstyle, numbers=left, firstnumber=268]
// Open-loop odometry integration
x_     += linear_  * std::cos(theta_) * dt;
y_     += linear_  * std::sin(theta_) * dt;
theta_ += angular_ * dt;
\end{lstlisting}
\end{filebox}

\textbf{Уравнения кинематики одноколейного робота}
(дискретная интеграция методом Эйлера вперёд):

\begin{align}
  x_{k+1} &= x_k + v \cos(\theta_k) \cdot \Delta t \label{eq:odom_x}\\
  y_{k+1} &= y_k + v \sin(\theta_k) \cdot \Delta t \label{eq:odom_y}\\
  \theta_{k+1} &= \theta_k + \omega \cdot \Delta t \label{eq:odom_theta}
\end{align}

Это модель \textbf{unicycle} (одноколейная модель): скорость $v$ направлена
вдоль текущего курса $\theta$.

\begin{notebox}
\textbf{Открытоконтурная одометрия}: модель не учитывает проскальзывание
гусениц. Именно для компенсации этой ошибки и используется Velocity EKF
(глава~\ref{ch:velocity_ekf}), который добавляет коррекцию от акселерометра.
\end{notebox}

\section{Кватернионное кодирование угла курса}

\begin{filebox}[title={\filepath{motor\_node.cpp}, строки 276--291 --- publishOdom()}]
\begin{lstlisting}[style=cppstyle, numbers=left, firstnumber=276]
void publishOdom(const rclcpp::Time & stamp)
{
  double sin_h = std::sin(theta_ / 2.0);
  double cos_h = std::cos(theta_ / 2.0);

  nav_msgs::msg::Odometry odom;
  odom.pose.pose.orientation.z = sin_h;
  odom.pose.pose.orientation.w = cos_h;
  odom.twist.twist.linear.x    = linear_;
  odom.twist.twist.angular.z   = angular_;
  odom_pub_->publish(odom);
}
\end{lstlisting}
\end{filebox}

Поворот в 3D вокруг вертикальной оси $Z$ кодируется кватернионом:

\begin{equation}
  \vect{q} = \bigl(\underbrace{0}_{q_w\text{'s missing part}},\,
  \underbrace{0}_{q_x},\,
  \underbrace{0}_{q_y},\,
  \underbrace{\sin(\theta/2)}_{q_z},\,
  \underbrace{\cos(\theta/2)}_{q_w}\bigr)
  \label{eq:quaternion_yaw}
\end{equation}

Для чистого вращения вокруг $Z$: $q_x = q_y = 0$, поэтому достаточно
хранить только $q_z$ и $q_w$ (что делает ROS2 сообщение \texttt{Odometry}).

\section{Сводная диаграмма}

\begin{center}
\begin{tikzpicture}[
  block/.style={rectangle, draw=blue!40, fill=blue!7, rounded corners,
                minimum width=3.5cm, minimum height=0.9cm,
                font=\small, align=center},
  arrow/.style={-{Stealth}, blue!60, thick}
]
  \node[block] (cmd) {Команда\\$(v_{\text{lin}},\, \omega_{\text{ang}})$};
  \node[block, right=2.5cm of cmd] (mix) {Микшер\\$L = \ell_{lin}+\ell_{ang}$\\$R = \ell_{lin}-\ell_{ang}$};
  \node[block, right=2.5cm of mix] (mot) {Моторы\\ШИМ $\in [-1, 1]$};
  \node[block, below=1.5cm of cmd] (odom) {Одометрия\\$x,y,\theta$};

  \draw[arrow] (cmd) -- (mix);
  \draw[arrow] (mix) -- (mot);
  \draw[arrow] (cmd) -- node[left, font=\tiny]{$v, \omega, \Delta t$} (odom);
\end{tikzpicture}
\end{center}

\section{Теоретическое обоснование: кинематика и динамика}

\subsection{Неголономная модель робота}

Гусеничный робот подчиняется \textbf{неголономному ограничению} ---
не может двигаться боком. В терминах ТАУ это выражается как:

\begin{equation}
  \dot{x}\sin\theta - \dot{y}\cos\theta = 0
  \label{eq:nonholonomic}
\end{equation}

Это ограничение не интегрируется в голономную связь вида $f(x, y, \theta) = 0$,
что делает систему \textbf{недоуправляемой} --- размерность конфигурационного
пространства $(x, y, \theta) = 3$ больше числа управлений $(v, \omega) = 2$.

\subsection{Модель unicycle в пространстве состояний}

Кинематическая модель~\eqref{eq:odom_x}--\eqref{eq:odom_theta} ---
это \textbf{нелинейная} система $\dot{\vect{x}} = \vect{f}(\vect{x}, \vect{u})$:

\begin{equation}
  \begin{pmatrix} \dot{x} \\ \dot{y} \\ \dot{\theta} \end{pmatrix}
  = \begin{pmatrix} \cos\theta & 0 \\ \sin\theta & 0 \\ 0 & 1 \end{pmatrix}
  \begin{pmatrix} v \\ \omega \end{pmatrix}
  \label{eq:unicycle_ss}
\end{equation}

Якобиан по состоянию (для EKF):

\begin{equation}
  \mat{F} = \frac{\partial \vect{f}}{\partial \vect{x}} =
  \begin{pmatrix}
    1 & 0 & -v\sin\theta\,\Delta t \\
    0 & 1 & v\cos\theta\,\Delta t \\
    0 & 0 & 1
  \end{pmatrix}
  \label{eq:unicycle_jacobian}
\end{equation}

\subsection{Связь с моделью двигателя}

Каждый борт робота приводится двигателем постоянного тока
(гл.~\ref{ch:theory}, раздел~\ref{sec:dc_motor}). Полная
динамическая модель включает электромеханику:

\begin{equation}
  T_a^2\ddot{\omega}_{\text{motor}} + T_y\dot{\omega}_{\text{motor}} + \omega_{\text{motor}} = k\,u_{\text{PWM}}
  \label{eq:motor_dynamics}
\end{equation}

Однако в проекте используется \textbf{кинематическая модель}
(без учёта инерции), что оправдано малым отношением
$T_a/T_{\text{control}} \ll 1$ (электрическая постоянная времени
много меньше периода управления $50\,\text{мс}$).

\subsection{Погрешность метода Эйлера}

Интегрирование методом Эйлера вперёд имеет погрешность $O(\Delta t^2)$
на шаг и $O(\Delta t)$ глобально. Для периода $\Delta t = 0.05\,\text{с}$
и $v_{\max} = 0.5\,\text{м/с}$ максимальная ошибка позиции за шаг:

\begin{equation}
  \epsilon \leq \frac{1}{2}|\ddot{x}|\Delta t^2 \approx \frac{v\omega\Delta t^2}{2}
  \leq \frac{0.5 \cdot 2.0 \cdot 0.0025}{2} = 1.25\,\text{мм}
  \label{eq:euler_error}
\end{equation}

Эта ошибка компенсируется Velocity EKF (гл.~\ref{ch:velocity_ekf}).

\section{Параметры}

\begin{center}
\begin{tabular}{lll}
\toprule
\textbf{Параметр} & \textbf{Значение} & \textbf{Описание}\\
\midrule
\texttt{max\_lin\_} & $0.5\,\text{м/с}$ & Максимальная линейная скорость\\
\texttt{max\_ang\_} & $2.0\,\text{рад/с}$ & Максимальная угловая скорость\\
Частота & $20\,\text{Гц}$ & Период управляющего цикла ($\Delta t = 0.05\,\text{с}$)\\
\bottomrule
\end{tabular}
\end{center}

% ============================================================
%  Глава 7 — Кватернионы и проекции сенсоров
%  Файл: compute_node/dashboard_node.py
% ============================================================
\chapter{Кватернионы, проекции и преобразования ориентации}
\label{ch:quaternions}

\begin{filebox}[title={Файл исходного кода}]
\filepath{compute\_node/dashboard\_node.py} \quad (Python, строки 298--521)
\end{filebox}

\section{Углы Эйлера из акселерометра (сырые)}

\begin{filebox}[title={\filepath{dashboard\_node.py}, строки 307--308 --- \_mqtt\_imu()}]
\begin{lstlisting}[style=pythonstyle, numbers=left, firstnumber=306]
# Raw pitch/roll from accel (no yaw)
pitch = math.degrees(math.atan2(ay, math.sqrt(ax*ax + az*az)))
roll  = math.degrees(math.atan2(-ax, az))
\end{lstlisting}
\end{filebox}

\textbf{Тангаж} (наклон вперёд/назад):

\begin{equation}
  \theta_{\text{pitch}} = \mathrm{atan2}\!\left(a_y,\, \sqrt{a_x^2 + a_z^2}\right)
  \label{eq:pitch_raw}
\end{equation}

\textbf{Крен} (наклон влево/вправо):

\begin{equation}
  \phi_{\text{roll}} = \mathrm{atan2}(-a_x,\, a_z)
  \label{eq:roll_raw}
\end{equation}

\begin{notebox}
Функция $\mathrm{atan2}(y, x)$ возвращает угол в $(-\pi, \pi]$ и
корректно обрабатывает все четыре квадранта, в отличие от $\arctan(y/x)$.
\end{notebox}

\section{Кватернион $\to$ угол курса (yaw из ROS2 Odometry)}

\begin{filebox}[title={\filepath{dashboard\_node.py}, строки 503--506 --- \_odom\_cb()}]
\begin{lstlisting}[style=pythonstyle, numbers=left, firstnumber=503]
q = msg.pose.pose.orientation
siny_cosp = 2.0 * (q.w * q.z + q.x * q.y)
cosy_cosp = 1.0 - 2.0 * (q.y * q.y + q.z * q.z)
yaw = math.atan2(siny_cosp, cosy_cosp)
\end{lstlisting}
\end{filebox}

\subsection{Кватернионы и ориентация в 3D}

Кватернион единичной длины $\vect{q} = (q_w, q_x, q_y, q_z)$, $|\vect{q}|=1$,
представляет вращение в 3D:

\begin{equation}
  \vect{q} = \cos\frac{\alpha}{2} + \hat{\vect{n}}\sin\frac{\alpha}{2}
  \label{eq:quaternion_def}
\end{equation}

где $\hat{\vect{n}}$ --- ось вращения, $\alpha$ --- угол.

\subsection{Формула извлечения угла курса}

Из матрицы вращения, построенной из кватерниона, угол вращения вокруг $Z$:

\begin{equation}
  R_{32} = 2(q_w q_z + q_x q_y), \quad
  R_{33} = 1 - 2(q_y^2 + q_z^2)
  \label{eq:rot_matrix_elems}
\end{equation}

\begin{equation}
  \boxed{\psi = \mathrm{atan2}\!\bigl(2(q_w q_z + q_x q_y),\;
    1 - 2(q_y^2 + q_z^2)\bigr)}
  \label{eq:yaw_from_quat}
\end{equation}

Это стандартная ZYX-декомпозиция (aerospace convention):
$\psi$ --- курс, $\theta$ --- тангаж, $\phi$ --- крен.

\section{Полярные координаты лазерного сканера $\to$ декартовы}

\begin{filebox}[title={\filepath{dashboard\_node.py}, строки 516--521 --- \_scan\_cb()}]
\begin{lstlisting}[style=pythonstyle, numbers=left, firstnumber=516]
def _scan_cb(self, msg: LaserScan):
    points, angle = [], msg.angle_min
    for r in msg.ranges:
        if msg.range_min < r < msg.range_max:
            points.append([
                round(r * math.cos(angle), 3),
                round(r * math.sin(angle), 3)
            ])
        angle += msg.angle_increment
\end{lstlisting}
\end{filebox}

\textbf{Полярные $\to$ декартовы координаты} (точка облака лазерного сканера):

\begin{equation}
  \begin{pmatrix} x \\ y \end{pmatrix}
  = r \begin{pmatrix} \cos\alpha \\ \sin\alpha \end{pmatrix}
  \label{eq:polar_to_cart}
\end{equation}

где $r$ --- измеренная дальность, $\alpha$ --- угол луча лазера.

\section{Геометрическое представление кватерниона}

\begin{center}
\begin{tikzpicture}[scale=1.2]
  % Coordinate axes
  \draw[-{Stealth}, thick] (0,0,0) -- (2,0,0) node[right]{$X$};
  \draw[-{Stealth}, thick] (0,0,0) -- (0,2,0) node[above]{$Z$};
  \draw[-{Stealth}, thick] (0,0,0) -- (0,0,2) node[right]{$Y$};

  % Rotation axis
  \draw[-{Stealth}, red!70, very thick] (0,0,0) -- (1.2, 1.5, 0.5)
    node[right, red!70]{$\hat{\vect{n}}$};

  % Arc for angle
  \draw[blue!60, thick, ->] (0.6,0,0) arc[start angle=0, end angle=50, radius=0.6]
    node[right, blue!60, font=\small]{$\alpha/2$};

  % Quaternion label
  \node[font=\small] at (0, -0.4, 0) {$\vect{q} = \cos\dfrac{\alpha}{2} + \hat{\vect{n}}\sin\dfrac{\alpha}{2}$};
\end{tikzpicture}
\end{center}

\section{Полная матрица вращения из кватерниона}

\begin{equation}
  \mat{R} = \begin{pmatrix}
    1-2(q_y^2+q_z^2) & 2(q_xq_y-q_wq_z) & 2(q_xq_z+q_wq_y) \\
    2(q_xq_y+q_wq_z) & 1-2(q_x^2+q_z^2) & 2(q_yq_z-q_wq_x) \\
    2(q_xq_z-q_wq_y) & 2(q_yq_z+q_wq_x) & 1-2(q_x^2+q_y^2)
  \end{pmatrix}
  \label{eq:rot_matrix_full}
\end{equation}

Для чистого вращения вокруг $Z$ ($q_x = q_y = 0$,
$q_z = \sin(\psi/2)$, $q_w = \cos(\psi/2)$):

\begin{equation}
  \mat{R}_z(\psi) = \begin{pmatrix}
    \cos\psi & -\sin\psi & 0 \\
    \sin\psi &  \cos\psi & 0 \\
    0 & 0 & 1
  \end{pmatrix}
  \label{eq:Rz_yaw}
\end{equation}

% ============================================================
%  Глава 8 — Калибровка IMU и цифровая фильтрация
%  Файл: pi_nodes/nodes/imu_node.py
% ============================================================
\chapter{Калибровка IMU и цифровая фильтрация}
\label{ch:imu_calibration}

\begin{filebox}[title={Файл исходного кода}]
\filepath{pi\_nodes/nodes/imu\_node.py} \quad (294 строки, Python)
\end{filebox}

\section{Аппаратная часть: MPU6050}

Датчик MPU6050 --- 6-осевая IMU (3-осевой акселерометр + 3-осевой гироскоп)
на шине I2C по адресу \texttt{0x68}.

\subsection{Масштабные коэффициенты}

\begin{filebox}[title={\filepath{imu\_node.py}, строки 46--55}]
\begin{lstlisting}[style=pythonstyle, numbers=left, firstnumber=46]
MPU6050_ADDR = 0x68
ACCEL_SCALE = 16384.0    # LSB/g at +/-2g
GYRO_SCALE  = 131.0      # LSB/(deg/s) at +/-250 deg/s
DEG2RAD = math.pi / 180.0
\end{lstlisting}
\end{filebox}

\textbf{Перевод сырых показаний в физические единицы:}

Акселерометр (из сырых 16-бит целых в м/с$^2$):

\begin{equation}
  a_{\text{м/с}^2} = \frac{\text{raw}_{16\text{bit}}}{16384} \times 9.81
  \label{eq:accel_scale}
\end{equation}

Гироскоп (из сырых 16-бит целых в рад/с):

\begin{equation}
  \omega_{\text{рад/с}} = \frac{\text{raw}_{16\text{bit}}}{131} \times \frac{\pi}{180}
  \label{eq:gyro_scale}
\end{equation}

\begin{filebox}[title={\filepath{imu\_node.py}, строки 133--149 --- \_read\_raw()}]
\begin{lstlisting}[style=pythonstyle, numbers=left, firstnumber=138]
raw = self._bus.read_i2c_block_data(MPU6050_ADDR, ACCEL_XOUT_H, 14)
ax = struct.unpack('>h', bytes(raw[0:2]))[0] / ACCEL_SCALE * 9.81
ay = struct.unpack('>h', bytes(raw[2:4]))[0] / ACCEL_SCALE * 9.81
az = struct.unpack('>h', bytes(raw[4:6]))[0] / ACCEL_SCALE * 9.81
gx = struct.unpack('>h', bytes(raw[8:10]))[0] / GYRO_SCALE * DEG2RAD
gy = struct.unpack('>h', bytes(raw[10:12]))[0] / GYRO_SCALE * DEG2RAD
gz = struct.unpack('>h', bytes(raw[12:14]))[0] / GYRO_SCALE * DEG2RAD
\end{lstlisting}
\end{filebox}

Считываются 14 байт начиная с регистра \texttt{0x3B}: три оси акселерометра (6 байт),
температура (2 байта, пропускается), три оси гироскопа (6 байт). Формат ---
big-endian signed 16-bit.

\section{Цифровой фильтр нижних частот (DLPF)}

\begin{filebox}[title={\filepath{imu\_node.py}, строки 87--89}]
\begin{lstlisting}[style=pythonstyle, numbers=left, firstnumber=87]
# DLPF mode 3: accel 44Hz, gyro 42Hz
bus.write_byte_data(MPU6050_ADDR, DLPF_CFG, 0x03)
bus.write_byte_data(MPU6050_ADDR, SMPLRT_DIV, 9)  # 1kHz/(1+9)=100Hz
\end{lstlisting}
\end{filebox}

Аппаратный DLPF на уровне сенсора: режим 3 обеспечивает полосу пропускания
44~Гц для акселерометра и 42~Гц для гироскопа, подавляя механические вибрации шасси.

Частота выборки:

\begin{equation}
  f_{\text{sample}} = \frac{f_{\text{osc}}}{1 + \text{SMPLRT\_DIV}} = \frac{1000}{1+9} = 100\,\text{Гц}
  \label{eq:sample_rate}
\end{equation}

\section{Экспоненциальное скользящее среднее (EMA)}

\begin{filebox}[title={\filepath{imu\_node.py}, строки 60--64, 151--167 --- \_apply\_ema()}]
\begin{lstlisting}[style=pythonstyle, numbers=left, firstnumber=60]
ACCEL_EMA_ALPHA = 0.2     # alpha for accel (8 Hz cutoff)
GYRO_EMA_ALPHA  = 0.5     # alpha for gyro  (less smoothing)
\end{lstlisting}
\begin{lstlisting}[style=pythonstyle, numbers=left, firstnumber=160]
self._ema_ax = a * ax + (1 - a) * self._ema_ax
self._ema_ay = a * ay + (1 - a) * self._ema_ay
self._ema_az = a * az + (1 - a) * self._ema_az
self._ema_gx = g * gx + (1 - g) * self._ema_gx
self._ema_gy = g * gy + (1 - g) * self._ema_gy
self._ema_gz = g * gz + (1 - g) * self._ema_gz
\end{lstlisting}
\end{filebox}

\textbf{Рекуррентная формула EMA} (фильтр первого порядка):

\begin{equation}
  \boxed{\hat{x}_k = \alpha \cdot x_k + (1-\alpha) \cdot \hat{x}_{k-1}}
  \label{eq:ema}
\end{equation}

Частота среза приблизительно:

\begin{equation}
  f_{\text{cutoff}} \approx \frac{\alpha \cdot f_s}{2\pi(1-\alpha)}
  \label{eq:ema_cutoff}
\end{equation}

Для акселерометра ($\alpha = 0.2$, $f_s = 50$~Гц):

\begin{equation}
  f_{\text{cutoff}} \approx \frac{0.2 \times 50}{2\pi \times 0.8} \approx 2.0\,\text{Гц}
  \label{eq:ema_cutoff_accel}
\end{equation}

\begin{center}
\begin{tabular}{llll}
\toprule
\textbf{Канал} & $\alpha$ & \textbf{Частота среза} & \textbf{Назначение}\\
\midrule
Акселерометр & $0.2$ & ${\approx}2\,\text{Гц}$ & Сильное подавление вибраций\\
Гироскоп     & $0.5$ & ${\approx}8\,\text{Гц}$ & Умеренное сглаживание\\
\bottomrule
\end{tabular}
\end{center}

\section{Калибровка при старте}

\begin{filebox}[title={\filepath{imu\_node.py}, строки 169--212 --- \_finish\_calibration()}]
\begin{lstlisting}[style=pythonstyle, numbers=left, firstnumber=171]
n = len(self._cal_samples)   # 100 samples (2 sec @ 50 Hz)
sum_ax = sum_ay = sum_az = 0.0
sum_gx = sum_gy = sum_gz = 0.0
for (ax, ay, az, gx, gy, gz) in self._cal_samples:
    sum_ax += ax; sum_ay += ay; sum_az += az
    sum_gx += gx; sum_gy += gy; sum_gz += gz

# Gyro offset: average at rest should be zero
self._gyro_offset = (sum_gx / n, sum_gy / n, sum_gz / n)

# Gravity vector in body frame at rest
avg_ax, avg_ay, avg_az = sum_ax/n, sum_ay/n, sum_az/n
self._gravity_body = (avg_ax, avg_ay, avg_az)

# Initial orientation from gravity
self._home_roll  = math.atan2(avg_ay, avg_az)
self._home_pitch = math.atan2(-avg_ax,
                    math.sqrt(avg_ay*avg_ay + avg_az*avg_az))
\end{lstlisting}
\end{filebox}

\subsection{Смещение нуля гироскопа}

В покое идеальный гироскоп показывает $\omega = 0$. В реальности
имеется постоянное смещение (bias):

\begin{equation}
  \vect{b}_{\text{gyro}} = \frac{1}{N}\sum_{k=1}^{N} \vect{\omega}_k^{\text{raw}}, \quad N = 100
  \label{eq:gyro_bias}
\end{equation}

Все последующие измерения корректируются:
$\vect{\omega}_{\text{corr}} = \vect{\omega}_{\text{raw}} - \vect{b}_{\text{gyro}}$.

\subsection{Вектор гравитации в системе координат тела}

Среднее показание акселерометра в покое --- это вектор гравитации
в системе координат IMU:

\begin{equation}
  \vect{g}_{\text{body}} = \frac{1}{N}\sum_{k=1}^{N} \vect{a}_k^{\text{raw}}
  = \bigl(\bar{a}_x,\, \bar{a}_y,\, \bar{a}_z\bigr)
  \label{eq:gravity_body}
\end{equation}

\subsection{Начальная ориентация (home)}

Из вектора гравитации вычисляется начальный наклон:

\textbf{Крен:}
\begin{equation}
  \phi_{\text{home}} = \mathrm{atan2}(\bar{a}_y,\, \bar{a}_z)
  \label{eq:home_roll}
\end{equation}

\textbf{Тангаж:}
\begin{equation}
  \theta_{\text{home}} = \mathrm{atan2}\!\left(-\bar{a}_x,\, \sqrt{\bar{a}_y^2 + \bar{a}_z^2}\right)
  \label{eq:home_pitch}
\end{equation}

Модуль вектора гравитации (для диагностики):

\begin{equation}
  |\vect{g}| = \sqrt{\bar{a}_x^2 + \bar{a}_y^2 + \bar{a}_z^2} \approx 9.81\,\text{м/с}^2
  \label{eq:g_magnitude}
\end{equation}

\section{Конвейер обработки IMU}

\begin{center}
\begin{tikzpicture}[
  node distance=0.5cm,
  block/.style={rectangle, draw=blue!50, fill=blue!8, rounded corners,
                minimum width=4cm, minimum height=0.7cm,
                font=\small\bfseries, align=center},
  arrow/.style={-{Stealth}, thick, blue!60}
]
  \node[block] (raw) {Сырые данные (14 байт I2C)};
  \node[block, below=of raw] (scale) {Масштабирование\\$a=\text{raw}/16384 \times 9.81$};
  \node[block, below=of scale] (ema) {EMA фильтр\\$\hat{x} = \alpha x + (1-\alpha)\hat{x}$};
  \node[block, below=of ema] (bias) {Коррекция смещения\\$\omega = \omega_{\text{raw}} - b$};
  \node[block, below=of bias] (ekf) {EKF IMU (глава~\ref{ch:imu_ekf})};
  \node[block, below=of ekf] (pub) {Публикация: $(\phi, \theta, \psi)$ + $\vect{a}$};

  \draw[arrow] (raw) -- (scale);
  \draw[arrow] (scale) -- (ema);
  \draw[arrow] (ema) -- (bias);
  \draw[arrow] (bias) -- (ekf);
  \draw[arrow] (ekf) -- (pub);
\end{tikzpicture}
\end{center}

\section{Сводная таблица параметров}

\begin{center}
\begin{tabular}{lll}
\toprule
\textbf{Параметр} & \textbf{Значение} & \textbf{Описание}\\
\midrule
\texttt{ACCEL\_SCALE} & $16384$ & LSB/g при $\pm 2\,g$\\
\texttt{GYRO\_SCALE}  & $131$   & LSB/(°/с) при $\pm 250\,\text{°/с}$\\
\texttt{DLPF\_CFG}    & $3$     & Полоса: 44~Гц акс., 42~Гц гиро\\
\texttt{SMPLRT\_DIV}  & $9$     & Частота: $100\,\text{Гц}$\\
\texttt{ACCEL\_EMA\_ALPHA} & $0.2$ & Сглаживание акселерометра\\
\texttt{GYRO\_EMA\_ALPHA}  & $0.5$ & Сглаживание гироскопа\\
\texttt{CALIBRATION\_SAMPLES} & $100$ & $2\,\text{сек}$ калибровки\\
\bottomrule
\end{tabular}
\end{center}

\section{C++ реализация (оптимизированная)}

\begin{filebox}[title={Файл исходного кода}]
\filepath{ros\_ws/src/robot\_pkg\_cpp/src/imu\_node.cpp} \quad (172 строки, C++)
\end{filebox}

\begin{filebox}[title={\filepath{imu\_node.cpp}, строки 90--127 --- readRaw()}]
\begin{lstlisting}[style=cppstyle, numbers=left, firstnumber=116]
// Bytes 0-5: accelerometer (X,Y,Z)
double ax = to_int16(raw[0],  raw[1])  / ACCEL_SCALE * GRAVITY;
double ay = to_int16(raw[2],  raw[3])  / ACCEL_SCALE * GRAVITY;
double az = to_int16(raw[4],  raw[5])  / ACCEL_SCALE * GRAVITY;
// Bytes 8-13: gyroscope (X,Y,Z)
double gx = to_int16(raw[8],  raw[9])  / GYRO_SCALE * DEG2RAD;
double gy = to_int16(raw[10], raw[11]) / GYRO_SCALE * DEG2RAD;
double gz = to_int16(raw[12], raw[13]) / GYRO_SCALE * DEG2RAD;
\end{lstlisting}
\end{filebox}

C++ версия использует прямой Linux I2C ioctl вместо Python smbus2.
Те же масштабные коэффициенты, тот же формат данных.
Преимущество: обратный вызов таймера ${\approx}10\,\mu\text{с}$ vs ${\approx}1\,\text{мс}$ в Python
(отсутствие GIL и интерпретатора).

\section{Теоретическое обоснование: цифровая фильтрация}

\subsection{EMA как дискретный ФНЧ}

Экспоненциальное скользящее среднее (EMA) --- это \textbf{дискретный
фильтр низких частот первого порядка} (IIR-фильтр):

\begin{equation}
  y_k = \alpha\, x_k + (1-\alpha)\, y_{k-1}
  \label{eq:ema_filter}
\end{equation}

$Z$-преобразование:

\begin{equation}
  H(z) = \frac{\alpha}{1 - (1-\alpha)z^{-1}}
  \label{eq:ema_ztransform}
\end{equation}

Полюс фильтра: $z_p = 1 - \alpha$. Условие устойчивости $|z_p| < 1$
выполняется при $0 < \alpha < 2$.

\subsection{Частотная характеристика EMA}

Частота среза (на уровне $-3\,\text{дБ}$):

\begin{equation}
  f_c = \frac{f_s}{2\pi}\arccos\!\left(\frac{2-\alpha}{2(1-\alpha)}\right)
  \label{eq:ema_cutoff}
\end{equation}

Для $\alpha = 0.15$, $f_s = 50\,\text{Гц}$: $f_c \approx 1.2\,\text{Гц}$.
Это подавляет вибрации двигателей ($> 10\,\text{Гц}$), сохраняя
полезный сигнал поворота ($< 1\,\text{Гц}$).

\subsection{Связь с аналоговым RC-фильтром}

EMA эквивалентен дискретизации аналогового RC-фильтра:

\begin{equation}
  W(p) = \frac{1}{\tau p + 1}, \quad \tau = \frac{(1-\alpha)\Delta t}{\alpha}
  \label{eq:ema_analog}
\end{equation}

Для $\alpha = 0.15$: $\tau = 0.85 \cdot 0.02 / 0.15 \approx 0.113\,\text{с}$.

\subsection{DLPF как аппаратный ФНЧ}

Цифровой фильтр низких частот (DLPF) MPU-6050 реализует
аналогичную функцию аппаратно, с полосой $44\,\text{Гц}$
(режим 3). Два последовательных ФНЧ (DLPF + EMA) дают
затухание $-12\,\text{дБ/окт}$ (второй порядок) в полосе
подавления.

% ============================================================
%  Глава 9 — A* Pathfinding и Bresenham Line-of-Sight
%  Файл: compute_node/simulator.py (строки 103--280)
% ============================================================
\chapter{A* Pathfinding и алгоритм Брезенхэма}
\label{ch:pathfinding}

\begin{filebox}[title={Файл исходного кода}]
\filepath{compute\_node/simulator.py} \quad (строки 103--280, Python)
\end{filebox}

\section{Постановка задачи}

Робот должен найти оптимальный путь от текущей позиции до цели,
обходя стены арены и запретные зоны. Используется \textbf{A*} ---
классический информированный алгоритм поиска на графе.

\section{Построение сетки занятости}

\begin{filebox}[title={\filepath{simulator.py}, строки 109--137 --- \_build\_grid()}]
\begin{lstlisting}[style=pythonstyle, numbers=left, firstnumber=109]
def _build_grid(arena, zones, robot_radius):
    """Build occupancy grid: True = blocked, False = free."""
    cols = int(arena.width / PATH_GRID_RES)   # 0.05 m/cell
    rows = int(arena.height / PATH_GRID_RES)
    grid = [[False] * cols for _ in range(rows)]
    margin = robot_radius + PATH_SAFETY_MARGIN  # 0.12 + 0.10 = 0.22 m

    for r in range(rows):
        for c in range(cols):
            wx = (c + 0.5) * PATH_GRID_RES
            wy = (r + 0.5) * PATH_GRID_RES
            # Block near arena walls
            if (wx < margin or wx > arena.width - margin or
                    wy < margin or wy > arena.height - margin):
                grid[r][c] = True
                continue
            # Block inside forbidden zones (with safety margin)
            for z in zones:
                zx1, zy1 = z['x1'] - margin, z['y1'] - margin
                zx2, zy2 = z['x2'] + margin, z['y2'] + margin
                if zx1 <= wx <= zx2 and zy1 <= wy <= zy2:
                    grid[r][c] = True
                    break
    return grid, rows, cols
\end{lstlisting}
\end{filebox}

Арена дискретизируется на ячейки $\delta = 5\,\text{см}$:

\begin{equation}
  \text{cols} = \left\lfloor \frac{W}{\delta} \right\rfloor, \quad
  \text{rows} = \left\lfloor \frac{H}{\delta} \right\rfloor
  \label{eq:grid_size}
\end{equation}

Зона безопасности вокруг каждого препятствия --- \textbf{раздутие Минковского}:

\begin{equation}
  \text{margin} = r_{\text{robot}} + r_{\text{safety}} = 0.12 + 0.10 = 0.22\,\text{м}
  \label{eq:minkowski_margin}
\end{equation}

\section{Преобразование координат}

\begin{filebox}[title={\filepath{simulator.py}, строки 140--147}]
\begin{lstlisting}[style=pythonstyle, numbers=left, firstnumber=140]
def _world_to_grid(wx, wy):
    return int(wx / PATH_GRID_RES), int(wy / PATH_GRID_RES)

def _grid_to_world(c, r):
    return (c + 0.5) * PATH_GRID_RES, (r + 0.5) * PATH_GRID_RES
\end{lstlisting}
\end{filebox}

\textbf{Мир $\to$ сетка:}
\begin{equation}
  c = \left\lfloor \frac{w_x}{\delta} \right\rfloor, \quad
  r = \left\lfloor \frac{w_y}{\delta} \right\rfloor
  \label{eq:world2grid_sim}
\end{equation}

\textbf{Сетка $\to$ мир} (центр ячейки):
\begin{equation}
  w_x = (c + 0.5)\,\delta, \quad w_y = (r + 0.5)\,\delta
  \label{eq:grid2world_sim}
\end{equation}

\section{Алгоритм A*}

\begin{filebox}[title={\filepath{simulator.py}, строки 150--213 --- find\_path()}]
\begin{lstlisting}[style=pythonstyle, numbers=left, firstnumber=173]
# A* with 8-directional movement
DIRS  = [(-1,0),(1,0),(0,-1),(0,1), (-1,-1),(-1,1),(1,-1),(1,1)]
COSTS = [1.0, 1.0, 1.0, 1.0,  1.414, 1.414, 1.414, 1.414]

def heuristic(r1, c1, r2, c2):
    dr = abs(r1 - r2)
    dc = abs(c1 - c2)
    return max(dr, dc) + 0.414 * min(dr, dc)  # octile distance

open_set = [(heuristic(sr, sc, gr, gc), 0.0, sr, sc)]
g_cost = {(sr, sc): 0.0}
came_from = {}

while open_set:
    _f, g, r, c = heapq.heappop(open_set)
    if r == gr and c == gc:
        # Reconstruct path ...
        break
    for (dr, dc), cost in zip(DIRS, COSTS):
        nr, nc = r + dr, c + dc
        if 0 <= nr < rows and 0 <= nc < cols and not grid[nr][nc]:
            ng = g + cost
            if ng < g_cost.get((nr, nc), float('inf')):
                g_cost[(nr, nc)] = ng
                f = ng + heuristic(nr, nc, gr, gc)
                came_from[(nr, nc)] = (r, c)
                heapq.heappush(open_set, (f, ng, nr, nc))
\end{lstlisting}
\end{filebox}

\subsection{Оценочная функция A*}

A* расширяет вершину с минимальным $f$:

\begin{equation}
  \boxed{f(n) = g(n) + h(n)}
  \label{eq:astar_f}
\end{equation}

где $g(n)$ --- стоимость пути от старта до $n$, $h(n)$ --- эвристика от $n$ до цели.

\subsection{Октильная эвристика}

Для 8-связного графа (с диагоналями) оптимальная допустимая эвристика:

\begin{equation}
  \boxed{h(n) = \max(\Delta r, \Delta c) + (\sqrt{2} - 1)\,\min(\Delta r, \Delta c)}
  \label{eq:octile}
\end{equation}

где $\Delta r = |r_n - r_g|$, $\Delta c = |c_n - c_g|$. В коде $\sqrt{2}-1 \approx 0.414$.

\textbf{Свойства:}
\begin{itemize}
  \item \textbf{Допустимая} (admissible): $h(n) \leq h^*(n)$ --- не переоценивает.
  \item \textbf{Монотонная} (consistent): $h(n) \leq c(n,n') + h(n')$ --- гарантирует
    оптимальность без повторного посещения.
\end{itemize}

\subsection{Стоимость переходов}

\begin{equation}
  c(n, n') = \begin{cases}
    1.0   & \text{кардинальное направление (4 стороны)} \\
    \sqrt{2} \approx 1.414 & \text{диагональное направление (4 угла)}
  \end{cases}
  \label{eq:move_cost}
\end{equation}

\section{Алгоритм Брезенхэма (Line-of-Sight)}

\begin{filebox}[title={\filepath{simulator.py}, строки 234--258 --- \_line\_of\_sight()}]
\begin{lstlisting}[style=pythonstyle, numbers=left, firstnumber=234]
def _line_of_sight(x0, y0, x1, y1, grid, rows, cols):
    """Bresenham: True if straight line is obstacle-free."""
    c0, r0 = _world_to_grid(x0, y0)
    c1, r1 = _world_to_grid(x1, y1)
    dc = abs(c1 - c0)
    dr = abs(r1 - r0)
    sc = 1 if c0 < c1 else -1
    sr = 1 if r0 < r1 else -1
    err = dc - dr
    while True:
        if grid[r0][c0]:
            return False         # blocked cell
        if r0 == r1 and c0 == c1:
            break
        e2 = 2 * err
        if e2 > -dr:
            err -= dr
            c0 += sc
        if e2 < dc:
            err += dc
            r0 += sr
    return True
\end{lstlisting}
\end{filebox}

\textbf{Алгоритм Брезенхэма} растеризует отрезок прямой на целочисленной сетке:

\begin{equation}
  \text{err} = \Delta c - \Delta r, \quad
  e_2 = 2 \cdot \text{err}
  \label{eq:bresenham_err}
\end{equation}

На каждом шаге:
\begin{itemize}
  \item Если $e_2 > -\Delta r$: $\text{err} \leftarrow \text{err} - \Delta r$, сдвиг по $c$.
  \item Если $e_2 < \Delta c$: $\text{err} \leftarrow \text{err} + \Delta c$, сдвиг по $r$.
\end{itemize}

Алгоритм работает за $O(\max(\Delta c, \Delta r))$ целочисленных операций (без деления и плавающей точки).

\section{Сглаживание пути (Line-of-Sight Pruning)}

\begin{filebox}[title={\filepath{simulator.py}, строки 261--279 --- \_smooth\_path()}]
\begin{lstlisting}[style=pythonstyle, numbers=left, firstnumber=261]
def _smooth_path(path, grid, rows, cols):
    """Reduce points via line-of-sight pruning."""
    smoothed = [path[0]]
    i = 0
    while i < len(path) - 1:
        best = i + 1
        for j in range(len(path) - 1, i + 1, -1):
            if _line_of_sight(path[i][0], path[i][1],
                              path[j][0], path[j][1],
                              grid, rows, cols):
                best = j
                break
        smoothed.append(path[best])
        i = best
    return smoothed
\end{lstlisting}
\end{filebox}

Удаляет промежуточные точки, если прямая видимость (Bresenham check) свободна.
Уменьшает число путевых точек с $O(N)$ до $O(\sqrt{N})$.

\section{Геометрия A* на сетке}

\begin{center}
\begin{tikzpicture}[scale=0.6]
  % Grid
  \draw[step=1, gray!30, very thin] (0,0) grid (8,6);

  % Obstacles
  \fill[red!20] (3,2) rectangle (5,4);
  \draw[red!60, thick] (3,2) rectangle (5,4);
  \node[font=\tiny, red!80] at (4,3) {Запретная зона};

  % Margin
  \draw[red!30, dashed] (2.5,1.5) rectangle (5.5,4.5);
  \node[font=\tiny, red!40] at (4,1.2) {margin};

  % Start
  \filldraw[blue!70] (1,1) circle (0.15);
  \node[below, font=\tiny, blue!70] at (1,0.85) {Start};

  % Goal
  \filldraw[green!70!black] (7,5) circle (0.15);
  \node[below, font=\tiny, green!70!black] at (7,4.85) {Goal};

  % Path
  \draw[blue!60, very thick, ->] (1,1) -- (2,2) -- (2,4.5) -- (3,5) -- (7,5);
  \node[font=\tiny, blue!60] at (0.5,3.5) {A* путь};
\end{tikzpicture}
\end{center}

\section{Сводная таблица параметров}

\begin{center}
\begin{tabular}{lll}
\toprule
\textbf{Параметр} & \textbf{Значение} & \textbf{Описание}\\
\midrule
\texttt{PATH\_GRID\_RES}     & $0.05\,\text{м}$ & Размер ячейки сетки\\
\texttt{PATH\_SAFETY\_MARGIN} & $0.10\,\text{м}$ & Запас безопасности\\
\texttt{ROBOT\_RADIUS}        & $0.12\,\text{м}$ & Радиус робота\\
Кардинальная стоимость        & $1.0$   & 4 направления\\
Диагональная стоимость        & $\sqrt{2} \approx 1.414$ & 4 диагонали\\
\bottomrule
\end{tabular}
\end{center}

% ============================================================
%  Глава 10 — Физический симулятор: кинематика, сенсоры, камера
%  Файл: compute_node/simulator.py
% ============================================================
\chapter{Физический симулятор: кинематика, сенсоры, камера}
\label{ch:simulator}

\begin{filebox}[title={Файл исходного кода}]
\filepath{compute\_node/simulator.py} \quad (строки 282--588, Python)
\end{filebox}

% ── Кинематика ─────────────────────────────────────────────────
\section{Кинематика робота (SimRobot)}

\begin{filebox}[title={\filepath{simulator.py}, строки 373--397 --- tick()}]
\begin{lstlisting}[style=pythonstyle, numbers=left, firstnumber=373]
def tick(self, dt: float):
    self._prev_v_linear = self.v_linear
    prev_x, prev_y = self.x, self.y
    # Integrate velocities (unicycle model)
    self.x += self.v_linear * math.cos(self.theta) * dt
    self.y += self.v_linear * math.sin(self.theta) * dt
    self.theta += self.v_angular * dt
    # Normalise theta
    self.theta = math.atan2(math.sin(self.theta), math.cos(self.theta))
    # Wall collision -- clamp inside arena
    margin = ROBOT_RADIUS
    self.x = max(margin, min(self.arena.width - margin, self.x))
    self.y = max(margin, min(self.arena.height - margin, self.y))
    # Forbidden zone collision -- push back
    for z in self.arena.forbidden_zones:
        zx1 = z['x1'] - margin
        zy1 = z['y1'] - margin
        zx2 = z['x2'] + margin
        zy2 = z['y2'] + margin
        if zx1 <= self.x <= zx2 and zy1 <= self.y <= zy2:
            self.x, self.y = prev_x, prev_y
            self.v_linear = 0.0
            break
\end{lstlisting}
\end{filebox}

\subsection{Дифференциальные уравнения кинематики}

Робот моделируется как \textbf{unicycle} --- одноколейная модель
с линейной и угловой скоростями:

\begin{align}
  \dot{x}      &= v \cos\theta \label{eq:sim_dx}\\
  \dot{y}      &= v \sin\theta \label{eq:sim_dy}\\
  \dot{\theta}  &= \omega       \label{eq:sim_dtheta}
\end{align}

Дискретизация методом Эйлера ($\Delta t = 0.05\,\text{с}$):

\begin{equation}
  \boxed{
  \begin{pmatrix} x_{k+1} \\ y_{k+1} \\ \theta_{k+1} \end{pmatrix}
  = \begin{pmatrix} x_k \\ y_k \\ \theta_k \end{pmatrix}
  + \Delta t \begin{pmatrix} v\cos\theta_k \\ v\sin\theta_k \\ \omega \end{pmatrix}
  }
  \label{eq:sim_euler}
\end{equation}

\subsection{Нормализация угла}

Для предотвращения накопления $\theta > 2\pi$:

\begin{equation}
  \theta \leftarrow \mathrm{atan2}\!\bigl(\sin\theta,\, \cos\theta\bigr)
  \label{eq:theta_norm}
\end{equation}

\subsection{Коллизия со стенами}

Ограничение позиции внутри арены:

\begin{equation}
  x \leftarrow \mathrm{clamp}(x,\; r_{\text{robot}},\; W - r_{\text{robot}}), \quad
  y \leftarrow \mathrm{clamp}(y,\; r_{\text{robot}},\; H - r_{\text{robot}})
  \label{eq:wall_collision}
\end{equation}

% ── Сенсоры ────────────────────────────────────────────────────
\section{Ультразвуковой сенсор (SimSensors)}

\begin{filebox}[title={\filepath{simulator.py}, строки 417--464 --- \_update\_ultrasonic()}]
\begin{lstlisting}[style=pythonstyle, numbers=left, firstnumber=417]
def _update_ultrasonic(self, robot, arena):
    """Ray-cast forward from robot."""
    rx, ry = robot.x, robot.y
    dx = math.cos(robot.theta)
    dy = math.sin(robot.theta)
    best = ULTRASONIC_MAX

    # Check walls (ray-plane intersection)
    if dx > 0:
        t = (arena.width - rx) / dx     # right wall
        if ULTRASONIC_MIN < t < best: best = t
    if dx < 0:
        t = -rx / dx                    # left wall
        if ULTRASONIC_MIN < t < best: best = t
    if dy > 0:
        t = (arena.height - ry) / dy    # top wall
        if ULTRASONIC_MIN < t < best: best = t
    if dy < 0:
        t = -ry / dy                    # bottom wall
        if ULTRASONIC_MIN < t < best: best = t

    # Check balls (ray-sphere intersection)
    for ball in arena.balls:
        if ball['grabbed']: continue
        to_ball_x = ball['x'] - rx
        to_ball_y = ball['y'] - ry
        proj = to_ball_x * dx + to_ball_y * dy   # projection
        perp = abs(to_ball_x * dy - to_ball_y * dx)  # perpendicular
        if proj > ULTRASONIC_MIN and perp < ball['radius'] + 0.05:
            if proj < best: best = proj

    # Gaussian noise
    noise = random.gauss(0, 0.005)
    self.range_m = max(ULTRASONIC_MIN, min(ULTRASONIC_MAX, best + noise))
\end{lstlisting}
\end{filebox}

\subsection{Пересечение луч--плоскость (стены)}

Луч из точки $(r_x, r_y)$ по направлению $(\cos\theta, \sin\theta)$:

\begin{equation}
  \vect{p}(t) = \begin{pmatrix} r_x \\ r_y \end{pmatrix}
  + t \begin{pmatrix} \cos\theta \\ \sin\theta \end{pmatrix}, \quad t > 0
  \label{eq:ray_def}
\end{equation}

Расстояние до стены $x = W$ (при $\cos\theta > 0$):

\begin{equation}
  t = \frac{W - r_x}{\cos\theta}
  \label{eq:ray_wall}
\end{equation}

\subsection{Пересечение луч--сфера (шарики)}

Для каждого шарика с центром $(b_x, b_y)$:

\textbf{Проекция} вектора ``к шарику'' на направление луча:

\begin{equation}
  t_{\text{proj}} = (\vect{b} - \vect{r}) \cdot \hat{\vect{d}}
  = (b_x - r_x)\cos\theta + (b_y - r_y)\sin\theta
  \label{eq:ball_proj}
\end{equation}

\textbf{Перпендикулярное расстояние} от луча до центра шарика:

\begin{equation}
  d_{\perp} = |(b_x - r_x)\sin\theta - (b_y - r_y)\cos\theta|
  \label{eq:ball_perp}
\end{equation}

Условие обнаружения: $d_\perp < r_{\text{ball}} + 0.05$ (конус ультразвука).

\subsection{Шумовая модель}

\begin{equation}
  r_{\text{meas}} = \mathrm{clamp}(r_{\text{true}} + \mathcal{N}(0, 0.005^2),\;
  r_{\min},\; r_{\max})
  \label{eq:ultrasonic_noise}
\end{equation}

% ── IMU ────────────────────────────────────────────────────────
\section{Симуляция IMU}

\begin{filebox}[title={\filepath{simulator.py}, строки 466--472 --- \_update\_imu()}]
\begin{lstlisting}[style=pythonstyle, numbers=left, firstnumber=466]
def _update_imu(self, robot):
    noise = random.gauss(0, 0.3)
    self.imu_yaw = math.degrees(robot.theta) + noise
    self.imu_pitch = random.gauss(0, 0.2)
    self.imu_roll  = random.gauss(0, 0.2)
    self.imu_gyro_z = robot.v_angular
    self.accel_x = (robot.v_linear - robot._prev_v_linear) / SIM_DT
\end{lstlisting}
\end{filebox}

Линейное ускорение из конечной разности скорости:

\begin{equation}
  a_x = \frac{v_k - v_{k-1}}{\Delta t}
  \label{eq:sim_accel}
\end{equation}

% ── Камера ─────────────────────────────────────────────────────
\section{Модель перспективной камеры (SimDetector)}

\begin{filebox}[title={\filepath{simulator.py}, строки 510--543 --- \_render\_camera()}]
\begin{lstlisting}[style=pythonstyle, numbers=left, firstnumber=514]
# Vector from robot to ball
dx = bx - robot.x
dy = by - robot.y
dist = math.sqrt(dx*dx + dy*dy)

# Angle to ball relative to robot heading
angle = math.atan2(dy, dx) - robot.theta
angle = math.atan2(math.sin(angle), math.cos(angle))

if abs(angle) > CAM_FOV / 2:
    continue   # outside field of view

# Project to screen
screen_x = int(CAM_W/2 + (angle / (CAM_FOV/2)) * (CAM_W/2))

# Apparent size (pinhole camera model)
apparent_px = int(FOCAL_LENGTH_PX * BALL_DIAMETER_M / dist)

# Vertical position (depth cue)
screen_y = horizon_y + int((1.0 - 0.03/max(0.1, dist)) *
                            (CAM_H - horizon_y) * 0.7)

# Confidence proportional to distance
conf = max(0.5, min(0.99, 1.0 - dist / 3.0))
\end{lstlisting}
\end{filebox}

\subsection{Проверка видимости (Field of View)}

Угол на объект относительно курса робота:

\begin{equation}
  \alpha = \mathrm{atan2}(b_y - r_y,\, b_x - r_x) - \theta_{\text{robot}}
  \label{eq:fov_angle}
\end{equation}

Объект виден если: $|\alpha| \leq \text{FOV}/2 = 30^\circ$.

\subsection{Горизонтальная проекция на экран}

\begin{equation}
  \boxed{s_x = \frac{W}{2} + \frac{\alpha}{\text{FOV}/2} \cdot \frac{W}{2}}
  \label{eq:screen_x}
\end{equation}

Это линейная проекция: угол $\alpha = 0$ --- центр кадра, $\pm \text{FOV}/2$ --- края.

\subsection{Модель pinhole-камеры (размер объекта)}

Видимый размер объекта на сенсоре:

\begin{equation}
  \boxed{s_{\text{px}} = \frac{f \cdot D_{\text{real}}}{d}}
  \label{eq:pinhole}
\end{equation}

где $f = 500\,\text{пкс}$ --- фокусное расстояние, $D_{\text{real}} = 0.04\,\text{м}$ ---
диаметр шарика, $d$ --- расстояние до объекта.

\begin{notebox}
Это обратная зависимость: при удвоении расстояния видимый размер уменьшается вдвое.
Тот же принцип используется в \texttt{follow\_me\_node.py}
(глава~\ref{ch:follow_me}) для оценки дистанции до человека.
\end{notebox}

\subsection{Вертикальная позиция (глубинная подсказка)}

Объекты на полу: чем ближе к роботу, тем ниже на экране:

\begin{equation}
  s_y = h_{\text{горизонт}} + \left(1 - \frac{0.03}{\max(0.1, d)}\right) \cdot 0.7(H - h_{\text{горизонт}})
  \label{eq:screen_y}
\end{equation}

\subsection{Модель доверия (confidence)}

\begin{equation}
  \text{conf} = \mathrm{clamp}\!\left(1 - \frac{d}{3.0},\; 0.5,\; 0.99\right)
  \label{eq:sim_conf}
\end{equation}

% ── Таблица ────────────────────────────────────────────────────
\section{Параметры симулятора}

\begin{center}
\begin{tabular}{lll}
\toprule
\textbf{Параметр} & \textbf{Значение} & \textbf{Описание}\\
\midrule
\texttt{ARENA\_W, H}        & $3.0 \times 3.0\,\text{м}$ & Размер арены\\
\texttt{SIM\_DT}            & $0.05\,\text{с}$ & Шаг интегрирования (20~Гц)\\
\texttt{MAX\_LINEAR}        & $0.3\,\text{м/с}$ & Макс. линейная скорость\\
\texttt{MAX\_ANGULAR}       & $2.0\,\text{рад/с}$ & Макс. угловая скорость\\
\texttt{ROBOT\_RADIUS}      & $0.12\,\text{м}$ & Радиус робота\\
\texttt{BALL\_DIAMETER\_M}  & $0.04\,\text{м}$ & Диаметр шарика\\
\texttt{FOCAL\_LENGTH\_PX}  & $500\,\text{пкс}$ & Фокусное расстояние камеры\\
\texttt{CAM\_FOV}           & $60^\circ$ & Поле зрения камеры\\
\texttt{CAM\_W, H}          & $640 \times 480\,\text{пкс}$ & Разрешение камеры\\
\texttt{ULTRASONIC\_MAX}    & $2.0\,\text{м}$ & Макс. дальность УЗ\\
\bottomrule
\end{tabular}
\end{center}

% ============================================================
%  Глава 11 — Следование за человеком: оценка дистанции и P-регулятор
%  Файл: ros_ws/src/robot_pkg/robot_pkg/follow_me_node.py
% ============================================================
\chapter{Следование за человеком: оценка дистанции и P-регулятор}
\label{ch:follow_me}

\begin{filebox}[title={Файл исходного кода}]
\filepath{ros\_ws/src/robot\_pkg/robot\_pkg/follow\_me\_node.py} \quad (138 строк, Python)
\end{filebox}

\section{Постановка задачи}

Модуль обнаруживает человека по детекции YOLO (класс ``person'') и
управляет роботом для поддержания заданной дистанции следования.

\section{Оценка расстояния по bounding box}

\begin{filebox}[title={\filepath{follow\_me\_node.py}, строки 83--88}]
\begin{lstlisting}[style=pythonstyle, numbers=left, firstnumber=83]
# Estimate distance from bounding box height
bbox_h = best.get('h', 1)
if bbox_h > 10:
    self._person_distance = (PERSON_HEIGHT_M * FOCAL_LENGTH_PX) / bbox_h
else:
    self._person_distance = TARGET_DISTANCE
\end{lstlisting}
\end{filebox}

Используется \textbf{модель pinhole-камеры} (та же, что в главе~\ref{ch:simulator}):

\begin{equation}
  h_{\text{px}} = \frac{f \cdot H_{\text{real}}}{d}
  \quad \Rightarrow \quad
  \boxed{d = \frac{H_{\text{person}} \cdot f}{h_{\text{bbox}}}}
  \label{eq:distance_estimation}
\end{equation}

где:
\begin{itemize}
  \item $H_{\text{person}} = 1.7\,\text{м}$ --- средний рост человека
  \item $f = 500\,\text{пкс}$ --- фокусное расстояние камеры
  \item $h_{\text{bbox}}$ --- высота bounding box в пикселях
\end{itemize}

\section{Рулевое управление (P-регулятор по углу)}

\begin{filebox}[title={\filepath{follow\_me\_node.py}, строки 90--103}]
\begin{lstlisting}[style=pythonstyle, numbers=left, firstnumber=90]
# Steering: centre person in frame
person_cx = best.get('x', IMG_CX) + best.get('w', 0) / 2.0
error_x = (person_cx - IMG_CX) / IMG_CX    # normalized -1 to +1

# Speed: proportional to distance error
dist_error = self._person_distance - TARGET_DISTANCE

twist = Twist()
twist.angular.z = -error_x * 0.8           # angular P-controller
if dist_error > 0.15:
    twist.linear.x = min(0.20, dist_error * 0.3)
elif dist_error < -0.15:
    twist.linear.x = max(-0.10, dist_error * 0.2)
\end{lstlisting}
\end{filebox}

\subsection{Нормализованная ошибка центрирования}

\begin{equation}
  e_x = \frac{c_{\text{person}} - c_{\text{image}}}{c_{\text{image}}} \in [-1, +1]
  \label{eq:center_error}
\end{equation}

где $c_{\text{image}} = 320\,\text{пкс}$ --- центр изображения.

\subsection{Угловая скорость (P-регулятор)}

\begin{equation}
  \omega = -K_\omega \cdot e_x, \quad K_\omega = 0.8
  \label{eq:angular_p}
\end{equation}

Знак <<$-$>> обеспечивает: человек справа ($e_x > 0$) $\Rightarrow$ робот поворачивает вправо ($\omega < 0$).

\subsection{Линейная скорость (P-регулятор с мёртвой зоной)}

Ошибка дистанции:

\begin{equation}
  e_d = d_{\text{measured}} - d_{\text{target}}, \quad d_{\text{target}} = 1.0\,\text{м}
  \label{eq:dist_error}
\end{equation}

\begin{equation}
  v = \begin{cases}
    \min(0.20,\; 0.3 \cdot e_d) & \text{если } e_d > 0.15\,\text{м (далеко)} \\
    \max(-0.10,\; 0.2 \cdot e_d) & \text{если } e_d < -0.15\,\text{м (близко)} \\
    0 & \text{иначе (мёртвая зона)}
  \end{cases}
  \label{eq:linear_p}
\end{equation}

Мёртвая зона $\pm 0.15\,\text{м}$ предотвращает колебания вокруг целевой дистанции.

\section{Теоретическое обоснование: замкнутая система слежения}

\subsection{Двухконтурная структура}

Система Follow-Me реализует \textbf{двухконтурный} регулятор
(гл.~\ref{ch:theory}, раздел~\ref{sec:p_controller}):

\begin{enumerate}
  \item \textbf{Угловой контур}: $\omega = -K_\omega e_x$ ---
    центрирование человека в кадре
  \item \textbf{Дистанционный контур}: $v = K_v e_d$ ---
    поддержание заданной дистанции
\end{enumerate}

\subsection{Передаточные функции контуров}

Для углового контура (робот --- интегратор $\dot{\theta} = \omega$):

\begin{equation}
  W_{\omega}(p) = \frac{K_\omega}{p + K_\omega} = \frac{0.8}{p + 0.8},
  \quad \tau_\omega = \frac{1}{K_\omega} = 1.25\,\text{с}
  \label{eq:follow_angular_tf}
\end{equation}

Для дистанционного контура (робот --- интегратор $\dot{d} = -v$):

\begin{equation}
  W_v(p) = \frac{K_v}{p + K_v} = \frac{0.3}{p + 0.3},
  \quad \tau_v = \frac{1}{K_v} = 3.3\,\text{с}
  \label{eq:follow_dist_tf}
\end{equation}

Угловой контур быстрее дистанционного ($\tau_\omega < \tau_v$) ---
сначала робот поворачивается к человеку, затем подъезжает.

\subsection{Мёртвая зона и устойчивость}

Мёртвая зона $\pm 0.15\,\text{м}$ предотвращает \textbf{автоколебания}
(предельный цикл), которые возникали бы из-за дискретности
и задержки обработки YOLO-детекции. Это аналог \textbf{гистерезиса}
в релейных системах управления.

\section{Таймаут потери цели}

\begin{filebox}[title={\filepath{follow\_me\_node.py}, строки 106--112}]
\begin{lstlisting}[style=pythonstyle, numbers=left, firstnumber=106]
def _tick(self):
    if self._tracking and (time.time() - self._last_person_time) > LOST_TIMEOUT:
        self._tracking = False
        self._cmd_pub.publish(Twist())  # stop
\end{lstlisting}
\end{filebox}

Если человек не обнаружен в течение $t_{\text{timeout}} = 2.0\,\text{с}$,
робот останавливается. Это защита от неконтролируемого движения.

\section{Сводная таблица параметров}

\begin{center}
\begin{tabular}{lll}
\toprule
\textbf{Параметр} & \textbf{Значение} & \textbf{Описание}\\
\midrule
\texttt{TARGET\_DISTANCE}  & $1.0\,\text{м}$ & Целевая дистанция следования\\
\texttt{PERSON\_HEIGHT\_M} & $1.7\,\text{м}$ & Средний рост человека\\
\texttt{FOCAL\_LENGTH\_PX} & $500\,\text{пкс}$ & Фокусное расстояние камеры\\
\texttt{IMG\_CX}           & $320\,\text{пкс}$ & Центр изображения по $x$\\
$K_\omega$                  & $0.8$ & Коэффициент P-регулятора угла\\
$K_v^+$                     & $0.3$ & Коэффициент скорости (далеко)\\
$K_v^-$                     & $0.2$ & Коэффициент скорости (близко)\\
\texttt{LOST\_TIMEOUT}     & $2.0\,\text{с}$ & Таймаут потери цели\\
\bottomrule
\end{tabular}
\end{center}

% ============================================================
%  Глава 12 — Патрулирование: навигация по точкам и кватернионы
%  Файл: ros_ws/src/robot_pkg/robot_pkg/patrol_node.py
% ============================================================
\chapter{Патрулирование: навигация по точкам и кватернионы}
\label{ch:patrol}

\begin{filebox}[title={Файл исходного кода}]
\filepath{ros\_ws/src/robot\_pkg/robot\_pkg/patrol\_node.py} \quad (145 строк, Python)
\end{filebox}

\section{Постановка задачи}

Модуль реализует циклический обход набора путевых точек (waypoints)
с отправкой целей в Nav2 (ROS2 навигационный стек).

\section{Проверка достижения точки}

\begin{filebox}[title={\filepath{patrol\_node.py}, строки 90--108 --- \_tick()}]
\begin{lstlisting}[style=pythonstyle, numbers=left, firstnumber=102]
# Check if reached
dist = math.hypot(tx - self._robot_x, ty - self._robot_y)
if dist < GOAL_TOLERANCE:   # 0.20 m
    self._current_idx = (self._current_idx + 1) % len(self._waypoints)
    self._goal_sent = False
\end{lstlisting}
\end{filebox}

\textbf{Евклидово расстояние} до цели:

\begin{equation}
  d = \sqrt{(t_x - x_r)^2 + (t_y - y_r)^2} = \mathrm{hypot}(\Delta x, \Delta y)
  \label{eq:patrol_dist}
\end{equation}

Точка считается достигнутой при $d < 0.20\,\text{м}$.

\textbf{Циклическая индексация:}

\begin{equation}
  i_{\text{next}} = (i_{\text{curr}} + 1) \bmod N_{\text{waypoints}}
  \label{eq:cyclic_idx}
\end{equation}

\section{Формирование цели Nav2 (кватернион из yaw)}

\begin{filebox}[title={\filepath{patrol\_node.py}, строки 110--118 --- \_send\_goal()}]
\begin{lstlisting}[style=pythonstyle, numbers=left, firstnumber=110]
def _send_goal(self, x, y, yaw):
    goal = PoseStamped()
    goal.header.stamp = self.get_clock().now().to_msg()
    goal.header.frame_id = 'map'
    goal.pose.position.x = x
    goal.pose.position.y = y
    goal.pose.orientation.z = math.sin(yaw / 2.0)
    goal.pose.orientation.w = math.cos(yaw / 2.0)
    self._goal_pub.publish(goal)
\end{lstlisting}
\end{filebox}

\textbf{Кватернион из угла курса} (чистое вращение вокруг $Z$):

\begin{equation}
  \boxed{
  q_z = \sin\!\left(\frac{\psi}{2}\right), \quad
  q_w = \cos\!\left(\frac{\psi}{2}\right), \quad
  q_x = q_y = 0
  }
  \label{eq:yaw_to_quat}
\end{equation}

Это следствие формулы Родрига: для оси вращения $\hat{\vect{n}} = (0, 0, 1)$
(ось $Z$, вертикальная):

\begin{equation}
  \vect{q} = \cos\frac{\psi}{2} + \hat{\vect{n}}\sin\frac{\psi}{2}
  = \left(0,\, 0,\, \sin\frac{\psi}{2},\, \cos\frac{\psi}{2}\right)
  \label{eq:rodriguez_z}
\end{equation}

\begin{notebox}
ROS2 использует порядок $(q_x, q_y, q_z, q_w)$ в сообщениях
\texttt{geometry\_msgs/Quaternion}. Для 2D-робота с вращением только
вокруг $Z$ достаточно задать $q_z$ и $q_w$, остальные равны нулю.
\end{notebox}

\section{Геометрия патрулирования}

\begin{center}
\begin{tikzpicture}[scale=1.5]
  % Waypoints
  \filldraw[green!70!black] (0, 0)   circle (0.06) node[below left, font=\tiny]{WP$_0$};
  \filldraw[green!70!black] (2, 0.5) circle (0.06) node[below right, font=\tiny]{WP$_1$};
  \filldraw[green!70!black] (2, 2)   circle (0.06) node[above right, font=\tiny]{WP$_2$};
  \filldraw[green!70!black] (0.5, 1.5) circle (0.06) node[above left, font=\tiny]{WP$_3$};

  % Path
  \draw[-{Stealth}, blue!60, thick] (0,0) -- (2,0.5);
  \draw[-{Stealth}, blue!60, thick] (2,0.5) -- (2,2);
  \draw[-{Stealth}, blue!60, thick] (2,2) -- (0.5,1.5);
  \draw[-{Stealth}, blue!60, thick, dashed] (0.5,1.5) -- (0,0);

  % Robot
  \filldraw[blue!80] (1.2, 0.3) circle (0.08);
  \node[below, font=\tiny, blue!80] at (1.2, 0.22) {Робот};

  % Distance
  \draw[red!50, dashed] (1.2, 0.3) -- (2, 0.5)
    node[midway, above, font=\tiny]{$d < 0.20\,\text{м}$?};

  % Tolerance circle
  \draw[green!50, dashed] (2, 0.5) circle (0.2);
\end{tikzpicture}
\end{center}

\section{Параметры}

\begin{center}
\begin{tabular}{lll}
\toprule
\textbf{Параметр} & \textbf{Значение} & \textbf{Описание}\\
\midrule
\texttt{GOAL\_TOLERANCE} & $0.20\,\text{м}$ & Допуск достижения точки\\
Частота обновления & $2\,\text{Гц}$ & Частота проверки цели\\
Кватернион & $(0, 0, \sin\frac{\psi}{2}, \cos\frac{\psi}{2})$ & Ориентация цели\\
\bottomrule
\end{tabular}
\end{center}

% ============================================================
%  Глава 13 — Геозона, преобразование сенсоров, реактивная навигация
%  Файлы: geofence_map_publisher.py, depth_to_scan_node.py, fallback_nav_node.py
% ============================================================
\chapter{Геозона, преобразование сенсоров и реактивная навигация}
\label{ch:geofence_sensors}

% ═══════════════════════════════════════════════════════════════
\section{Геозона: алгоритм Point-in-Polygon}

\begin{filebox}[title={Файл исходного кода}]
\filepath{compute\_node/geofence\_map\_publisher.py} \quad (108 строк, Python)
\end{filebox}

\begin{filebox}[title={\filepath{geofence\_map\_publisher.py}, строки 79--91 --- \_point\_in\_polygon()}]
\begin{lstlisting}[style=pythonstyle, numbers=left, firstnumber=79]
@staticmethod
def _point_in_polygon(x, y, poly):
    """Ray-casting algorithm."""
    n = len(poly)
    inside = False
    j = n - 1
    for i in range(n):
        xi, yi = poly[i]
        xj, yj = poly[j]
        if ((yi > y) != (yj > y)) and \
           (x < (xj - xi) * (y - yi) / (yj - yi) + xi):
            inside = not inside
        j = i
    return inside
\end{lstlisting}
\end{filebox}

\subsection{Алгоритм трассировки луча (Ray Casting)}

Для определения, находится ли точка $(x, y)$ внутри полигона,
из неё проводится горизонтальный луч вправо и считается число
пересечений с рёбрами:

\begin{equation}
  \text{inside} = \bigl(\text{число пересечений}\bigr) \bmod 2 = 1
  \label{eq:raycasting}
\end{equation}

Для каждого ребра $(x_i, y_i) \to (x_j, y_j)$:

\textbf{1. Проверка пересечения по $y$:}
\begin{equation}
  (y_i > y) \neq (y_j > y) \quad\text{(ребро пересекает горизонталь $y$)}
  \label{eq:edge_cross}
\end{equation}

\textbf{2. Вычисление $x$-координаты пересечения} (линейная интерполяция):
\begin{equation}
  x_{\text{cross}} = x_i + \frac{(x_j - x_i)(y - y_i)}{y_j - y_i}
  \label{eq:x_intersect}
\end{equation}

\textbf{3. Если} $x < x_{\text{cross}}$ --- луч пересекает ребро, инвертируем флаг.

\subsection{Построение карты занятости}

\begin{filebox}[title={\filepath{geofence\_map\_publisher.py}, строки 46--77}]
\begin{lstlisting}[style=pythonstyle, numbers=left, firstnumber=47]
poly = np.array(self._polygon, dtype=np.float64)
x_min = poly[:, 0].min() - self._padding   # 0.5 m padding
x_max = poly[:, 0].max() + self._padding
width  = int((x_max - x_min) / self._resolution)  # 0.05 m/cell
height = int((y_max - y_min) / self._resolution)

# 0=free inside, 100=lethal outside
grid = np.full((height, width), 100, dtype=np.int8)
for iy in range(height):
    for ix in range(width):
        px = x_min + ix * self._resolution
        py = y_min + iy * self._resolution
        if self._point_in_polygon(px, py, poly):
            grid[iy, ix] = 0
\end{lstlisting}
\end{filebox}

Карта создаётся как \texttt{OccupancyGrid} (ROS2 nav\_msgs):
\begin{itemize}
  \item $0$ --- свободная ячейка (внутри полигона)
  \item $100$ --- летальное препятствие (за пределами арены)
\end{itemize}

% ═══════════════════════════════════════════════════════════════
\section{Ультразвук $\to$ LaserScan: синтетическое сканирование}

\begin{filebox}[title={Файл исходного кода}]
\filepath{compute\_node/depth\_to\_scan\_node.py} \quad (82 строки, Python)
\end{filebox}

\begin{filebox}[title={\filepath{depth\_to\_scan\_node.py}, строки 43--65}]
\begin{lstlisting}[style=pythonstyle, numbers=left, firstnumber=43]
def _publish_scan(self):
    scan = LaserScan()
    scan.angle_min = -math.pi / 2.0        # -90 deg
    scan.angle_max =  math.pi / 2.0        # +90 deg
    scan.angle_increment = (ANGLE_MAX - ANGLE_MIN) / NUM_BEAMS  # 180 beams
    scan.range_min = 0.02
    scan.range_max = 3.0

    ranges = [MAX_RANGE] * NUM_BEAMS        # all unknown
    centre = NUM_BEAMS // 2
    # Spread ultrasonic over ~15 deg cone
    cone_beams = int(0.26 / scan.angle_increment)
    for i in range(max(0, centre - cone_beams),
                   min(NUM_BEAMS, centre + cone_beams + 1)):
        ranges[i] = min(self._latest_range, MAX_RANGE)
    scan.ranges = ranges
\end{lstlisting}
\end{filebox}

\subsection{Геометрия конуса ультразвука}

Один ультразвуковой датчик HC-SR04 даёт единственное расстояние
с конусом обзора ${\approx}30^\circ$. Для совместимости с slam\_toolbox
это значение «размазывается» по нескольким лучам:

\begin{equation}
  \Delta\alpha = \frac{\alpha_{\max} - \alpha_{\min}}{N_{\text{beams}}}
  = \frac{\pi}{180} \approx 0.0175\,\text{рад/луч}
  \label{eq:angle_incr}
\end{equation}

\begin{equation}
  N_{\text{cone}} = \left\lfloor \frac{0.26}{\Delta\alpha} \right\rfloor
  \approx 15\,\text{лучей в полуконусе}
  \label{eq:cone_beams}
\end{equation}

Итого ${\approx}30$ центральных лучей получают реальное измерение,
остальные 150 --- значение $r_{\max} = 3.0\,\text{м}$ (неизвестно).

% ═══════════════════════════════════════════════════════════════
\section{Реактивная навигация (fallback)}

\begin{filebox}[title={Файл исходного кода}]
\filepath{pi\_nodes/nodes/fallback\_nav\_node.py} \quad (176 строк, Python)
\end{filebox}

\begin{filebox}[title={\filepath{fallback\_nav\_node.py}, строки 122--150 --- \_run\_fallback\_drive()}]
\begin{lstlisting}[style=pythonstyle, numbers=left, firstnumber=136]
if r > SAFE_M:                          # r > 0.40 m
    cmd['linear_x'] = FWD_SPEED        # 0.10 m/s
    self._state = STATE_FORWARD
elif r > STOP_M:                        # 0.20 < r < 0.40
    scale = (r - STOP_M) / (SAFE_M - STOP_M)
    cmd['linear_x'] = SLOW_SPEED + scale * (FWD_SPEED - SLOW_SPEED)
    self._state = STATE_CAUTION
else:                                   # r < 0.20 m
    self._state = STATE_ROTATE
    cmd['angular_z'] = ROT_SPEED * self._rotate_dir  # +/- 0.8 rad/s
\end{lstlisting}
\end{filebox}

\subsection{Пропорциональное масштабирование скорости}

В зоне осторожности ($r_{\text{stop}} < r < r_{\text{safe}}$):

\begin{equation}
  \text{scale} = \frac{r - r_{\text{stop}}}{r_{\text{safe}} - r_{\text{stop}}}
  = \frac{r - 0.20}{0.40 - 0.20} \in [0, 1]
  \label{eq:caution_scale}
\end{equation}

\begin{equation}
  \boxed{v = v_{\text{slow}} + \text{scale} \cdot (v_{\text{fwd}} - v_{\text{slow}})}
  \label{eq:caution_speed}
\end{equation}

Это \textbf{линейная интерполяция} между $v_{\text{slow}} = 0.05\,\text{м/с}$
и $v_{\text{fwd}} = 0.10\,\text{м/с}$.

\subsection{Конечный автомат реактивной навигации}

\begin{center}
\begin{tikzpicture}[
  state/.style={rectangle, rounded corners, draw=blue!60, fill=blue!10,
                minimum width=2.5cm, minimum height=0.9cm, font=\small\bfseries},
  arrow/.style={-{Stealth[length=6pt]}, thick}
]
  \node[state] (F) {FORWARD};
  \node[state, right=3cm of F] (C) {CAUTION};
  \node[state, below=1.5cm of C] (R) {ROTATE};

  \draw[arrow] (F) to node[above, font=\tiny]{$r < 0.40$} (C);
  \draw[arrow] (C) to node[right, font=\tiny]{$r < 0.20$} (R);
  \draw[arrow] (R) to[bend right=30] node[left, font=\tiny]{$r > 0.40$} (F);
  \draw[arrow] (C) to[bend left=20] node[above, font=\tiny]{$r > 0.40$} (F);
\end{tikzpicture}
\end{center}

\subsection{Теоретическое обоснование: конечный автомат как дискретная система}

Реактивная навигация реализует \textbf{гибридную систему} ---
комбинацию непрерывной динамики (скорость) и дискретных переключений
(состояния автомата). В терминах ТАУ это \textbf{переключающаяся система}:

\begin{equation}
  \dot{\vect{x}} = \vect{f}_{q(t)}(\vect{x}), \quad q \in \{\text{FWD}, \text{CAUT}, \text{ROT}\}
  \label{eq:switching_system}
\end{equation}

\textbf{Устойчивость} переключающейся системы гарантирована, если
каждая подсистема устойчива и переключения не слишком частые
(условие dwell time).

Линейная интерполяция~\eqref{eq:caution_speed} в зоне CAUTION
обеспечивает \textbf{непрерывность} скорости на границах переключений,
предотвращая рывки.

\section{Сводная таблица параметров}

\begin{center}
\begin{tabular}{llll}
\toprule
\textbf{Модуль} & \textbf{Параметр} & \textbf{Значение} & \textbf{Описание}\\
\midrule
Геозона & Разрешение & $0.05\,\text{м}$ & Размер ячейки\\
Геозона & Отступ & $0.5\,\text{м}$ & Запас за полигоном\\
LaserScan & $N_{\text{beams}}$ & $180$ & Кол-во лучей\\
LaserScan & FOV & $180^\circ$ & Поле зрения\\
Fallback & $r_{\text{safe}}$ & $0.40\,\text{м}$ & Безопасная дистанция\\
Fallback & $r_{\text{stop}}$ & $0.20\,\text{м}$ & Стоп-дистанция\\
Fallback & $\omega_{\text{rot}}$ & $0.8\,\text{рад/с}$ & Скорость поворота\\
\bottomrule
\end{tabular}
\end{center}

% ============================================================
%  Глава 14 — ШИМ-управление моторами и PCA9685
%  Файл: ros_ws/src/robot_pkg_cpp/src/motor_node.cpp
% ============================================================
\chapter{ШИМ-управление моторами и PCA9685}
\label{ch:pwm_motor}

\begin{filebox}[title={Файл исходного кода}]
\filepath{ros\_ws/src/robot\_pkg\_cpp/src/motor\_node.cpp} \quad (343 строки, C++)
\end{filebox}

\section{Контроллер PCA9685}

PCA9685 --- 16-канальный ШИМ-контроллер на I2C (адрес \texttt{0x5F}).
Каждый канал формирует ШИМ-сигнал с 12-битным разрешением ($0\dots4095$).

\subsection{Вычисление прескейлера}

\begin{filebox}[title={\filepath{motor\_node.cpp}, строки 91--94}]
\begin{lstlisting}[style=cppstyle, numbers=left, firstnumber=91]
// prescale = round(osc / (4096 * freq)) - 1
uint8_t prescale = static_cast<uint8_t>(
  std::round(static_cast<double>(OSC_FREQ_HZ) /
             (4096.0 * static_cast<double>(PWM_FREQ_HZ))) - 1.0);
\end{lstlisting}
\end{filebox}

Формула прескейлера для заданной частоты ШИМ:

\begin{equation}
  \boxed{\text{prescale} = \left\lfloor \frac{f_{\text{osc}}}{4096 \cdot f_{\text{PWM}}} + 0.5 \right\rfloor - 1}
  \label{eq:prescale}
\end{equation}

Для $f_{\text{osc}} = 25\,\text{МГц}$, $f_{\text{PWM}} = 50\,\text{Гц}$:

\begin{equation}
  \text{prescale} = \left\lfloor \frac{25{,}000{,}000}{4096 \times 50} + 0.5 \right\rfloor - 1
  = \left\lfloor 122.07 + 0.5 \right\rfloor - 1 = 121
  \label{eq:prescale_val}
\end{equation}

\section{Slow-Decay режим управления мотором}

\begin{filebox}[title={\filepath{motor\_node.cpp}, строки 210--232 --- setMotor()}]
\begin{lstlisting}[style=cppstyle, numbers=left, firstnumber=210]
void setMotor(int idx, double throttle)
{
  throttle = std::clamp(throttle, -1.0, 1.0);
  int in1, in2;
  if (throttle > 0.001) {
    // Forward: IN1 = 4095,  IN2 = 4095*(1-t)
    in1 = PWM_MAX;
    in2 = static_cast<int>(PWM_MAX * (1.0 - throttle));
  } else if (throttle < -0.001) {
    // Reverse: IN1 = 4095*(1+t),  IN2 = 4095
    in1 = static_cast<int>(PWM_MAX * (1.0 + throttle));
    in2 = PWM_MAX;
  } else {
    // Active brake: IN1 = IN2 = 4095
    in1 = PWM_MAX;
    in2 = PWM_MAX;
  }
  pca_->setPwm(MOTORS[idx].in1, in1);
  pca_->setPwm(MOTORS[idx].in2, in2);
}
\end{lstlisting}
\end{filebox}

\subsection{Формулы Slow-Decay}

В режиме \textbf{Slow Decay} (медленное затухание тока) оба вывода
активны, что обеспечивает плавное торможение:

\textbf{Вперёд} ($t > 0$, throttle $\in (0, 1]$):
\begin{equation}
  \text{IN1} = 4095, \quad \text{IN2} = 4095 \cdot (1 - t)
  \label{eq:forward_pwm}
\end{equation}

\textbf{Назад} ($t < 0$, throttle $\in [-1, 0)$):
\begin{equation}
  \text{IN1} = 4095 \cdot (1 + t), \quad \text{IN2} = 4095
  \label{eq:reverse_pwm}
\end{equation}

\textbf{Активное торможение} ($t = 0$):
\begin{equation}
  \text{IN1} = \text{IN2} = 4095 \quad \text{(обе обмотки замкнуты)}
  \label{eq:brake_pwm}
\end{equation}

\subsection{Схема каналов моторов}

\begin{center}
\begin{tabular}{llll}
\toprule
\textbf{Мотор} & \textbf{Позиция} & \textbf{IN1 канал} & \textbf{IN2 канал}\\
\midrule
M1 & Правый передний & ch15 & ch14\\
M2 & Левый передний  & ch12 & ch13\\
M3 & Левый задний    & ch11 & ch10\\
M4 & Правый задний   & ch8  & ch9\\
\bottomrule
\end{tabular}
\end{center}

\section{Дифференциальный микшер}

\begin{filebox}[title={\filepath{motor\_node.cpp}, строки 251--266}]
\begin{lstlisting}[style=cppstyle, numbers=left, firstnumber=251]
double lin_t = (max_lin_ > 0.0) ? (linear_  / max_lin_) : 0.0;
double ang_t = (max_ang_ > 0.0) ? (angular_ / max_ang_) : 0.0;

double left  = lin_t + ang_t;
double right = lin_t - ang_t;

double max_val = std::max({std::abs(left), std::abs(right), 1.0});
if (max_val > 1.0) { left /= max_val; right /= max_val; }
\end{lstlisting}
\end{filebox}

Нормализованные команды ($\ell \in [-1, 1]$):

\begin{align}
  \ell_{\text{lin}} &= \frac{v}{v_{\max}}, \quad
  \ell_{\text{ang}} = \frac{\omega}{\omega_{\max}} \label{eq:norm_cmd}\\[6pt]
  \text{left}  &= \ell_{\text{lin}} + \ell_{\text{ang}} \label{eq:left_mix}\\
  \text{right} &= \ell_{\text{lin}} - \ell_{\text{ang}} \label{eq:right_mix}
\end{align}

\textbf{Нормализация} (сохранение пропорции при насыщении):

\begin{equation}
  m = \max(1, |\text{left}|, |\text{right}|), \quad
  \text{left} \leftarrow \frac{\text{left}}{m}, \quad
  \text{right} \leftarrow \frac{\text{right}}{m}
  \label{eq:mix_normalize}
\end{equation}

\section{Теоретическое обоснование: электропривод}

\subsection{Модель двигателя постоянного тока}

Электродвигатели робота описываются дифференциальным уравнением
(гл.~\ref{ch:theory}, раздел~\ref{sec:dc_motor}):

\begin{equation}
  L_a \frac{dI}{dt} + R_a I = U - C_e \omega, \qquad
  J \frac{d\omega}{dt} = C_m I - M_{\text{нагр}}
  \label{eq:motor_electrical}
\end{equation}

где $L_a$, $R_a$ --- индуктивность и сопротивление якоря, $I$ --- ток,
$C_e$ --- коэффициент ЭДС, $C_m$ --- коэффициент момента,
$J$ --- момент инерции, $M_{\text{нагр}}$ --- момент нагрузки.

Исключая ток, получаем передаточную функцию:

\begin{equation}
  W(p) = \frac{\omega(p)}{U(p)} = \frac{K_m}{T_m T_e p^2 + T_m p + 1}
  \label{eq:motor_full_tf}
\end{equation}

\subsection{Slow-Decay и эффективный момент}

В режиме Slow Decay эффективное напряжение на двигателе:

\begin{equation}
  U_{\text{эфф}} = U_{\text{пит}} \cdot t, \quad t \in [-1, 1]
  \label{eq:slow_decay_voltage}
\end{equation}

При активном торможении ($t = 0$, $\text{IN1} = \text{IN2} = 4095$)
обмотки замкнуты, и ЭДС двигателя создаёт тормозной ток:

\begin{equation}
  I_{\text{торм}} = \frac{C_e \omega}{R_a}, \qquad
  M_{\text{торм}} = C_m I_{\text{торм}} = \frac{C_e C_m}{R_a}\omega
  \label{eq:brake_torque}
\end{equation}

Это обеспечивает плавное торможение, пропорциональное скорости
(вязкое демпфирование).

\subsection{Замкнутый контур <<микшер + двигатель>>}

Передаточная функция от команды $(v_{\text{lin}}, \omega_{\text{ang}})$
до скоростей бортов:

\begin{equation}
  \begin{pmatrix} V_L(p) \\ V_R(p) \end{pmatrix}
  = \frac{K_m}{T_m p + 1}
  \begin{pmatrix} 1 & 1 \\ 1 & -1 \end{pmatrix}
  \begin{pmatrix} \ell_{\text{lin}}(p) \\ \ell_{\text{ang}}(p) \end{pmatrix}
  \label{eq:mixer_motor}
\end{equation}

При $T_m \ll T_{\text{control}}$ система упрощается до
статической модели $V_L \approx K_m(\ell_{\text{lin}} + \ell_{\text{ang}})$.

\section{Параметры}

\begin{center}
\begin{tabular}{lll}
\toprule
\textbf{Параметр} & \textbf{Значение} & \textbf{Описание}\\
\midrule
\texttt{PCA9685\_ADDR} & \texttt{0x5F} & I2C-адрес контроллера\\
\texttt{PWM\_FREQ\_HZ} & $50\,\text{Гц}$ & Частота ШИМ\\
\texttt{PWM\_MAX}       & $4095$ & 12-бит макс. значение\\
\texttt{WHEEL\_BASE}    & $0.17\,\text{м}$ & Расстояние между гусеницами\\
\bottomrule
\end{tabular}
\end{center}


\end{document}
